\section*{3. Théorèmes de finitude généraux}
\addcontentsline{toc}{section}{3. Théorèmes de finitude généraux}

\begin{lem}\label{lem:3.1} Soient $S$ un schéma et $X$ un $S$-schéma quasi-compact et quasi-séparé ; soit $\underline{\operatorname{Pic}}_{X/S}$ le foncteur de Picard (localisé pour la topologie f.p.p.f.)\textup{($*$)}. Alors pour tout système projectif filtrant $(Z_{\alpha})$ de $S$-schémas quasi-compacts et quasi-séparés, à morphismes de transition affines, l'application canonique
\[
\text{(3.1.1)} \qquad \varinjlim \underline{\operatorname{Pic}}_{X/S}(Z_{\alpha}) \to \underline{\operatorname{Pic}}_{X/S}(\varprojlim Z_{\alpha})
\]
est bijective.

\smallskip
\noindent($*$) Cf.\@ SGA 3 IV 6.3, et FGA n$^\circ$ 232 pour la définition de $\underline{\operatorname{Pic}}_{X/S}$, où il convient ici, comme indiqué, de remplacer la topologie fpqc par la topologie fppf.
\end{lem}

\noindent\textit{Par conséquent} (EGA IV 8.14.2), si le foncteur de Picard est représenté par un $S$-schéma $\underline{\operatorname{Pic}}_{X/S}$, alors $\underline{\operatorname{Pic}}_{X/S}$ est localement de présentation finie.

\noindent\textit{En effet,} un élément $\lambda_{\alpha}$ de $\underline{\operatorname{Pic}}_{X/S}(Z_{\alpha})$ est représenté par un faisceau inversible $L'_{\alpha}$ sur un $X_{Z'_{\alpha}}$ où $Z'_{\alpha} \to Z_{\alpha}$ est un morphisme f.p.p.f. Posons $Z = \varprojlim Z_{\alpha}$. L'image $\lambda$ de $\lambda_{\alpha}$ dans $\underline{\operatorname{Pic}}_{X/S}(Z)$ est représentée par le faisceau inversible $L'$ sur $X_{Z'}$, image inverse de $L'_{\alpha}$, où $Z' = Z'_{\alpha} \times_{Z_{\alpha}} Z$.

Si $\lambda$ est l'élément neutre de $\underline{\operatorname{Pic}}_{X/S}(Z)$, il existe un morphisme f.p.p.f.\@ $Z'' \to Z'$ tel que l'image inverse de $L'$ sur $X_{Z''}$ soit $\simeq \cO_{X_{Z''}}$.

% ============================================================
% PAGE 635
% ============================================================
\begin{flushright}
$\fbox{\text{XIII}}$
\end{flushright}

\noindent
En vertu de (EGA IV 8.8.2, 8.10.5 (vi), 11.2.6 (ii), 8.5.2 (i), et 8.5.2.4), il existe un $\beta \geq \alpha$ et un morphisme f.p.p.f.\@ $Z''_{\beta} \to Z'_{\beta}$ où $Z'_{\beta} = Z'_{\alpha} \times_{Z_{\alpha}} Z_{\beta}$ tel que l'image inverse de $L'_{\alpha}$ sur $X_{Z''_{\beta}}$ soit $\cO_{X_{Z''_{\beta}}}$.

L'image de $\lambda_{\alpha}$ est donc déjà l'élément neutre dans $\underline{\operatorname{Pic}}_{X/S}(Z_{\beta})$. Par suite 3.1.1 est injective.

La surjectivité de 3.1.1 se démontre d'une façon analogue.

\begin{prop}\label{prop:3.2} Soient $S$ un schéma noethérien et $X$ un $S$-schéma propre. On suppose que le foncteur de Picard est représenté par un $S$-schéma $\underline{\operatorname{Pic}}_{X/S}$. Alors :
\begin{enumerate}[label={\roman*)}]
\item $\underline{\operatorname{Pic}}_{X/S}$ est localement de type fini sur $S$.
\item Les parties $\Lambda$ de $\underline{\operatorname{Pic}}_{X/S}$ correspondent biunivoquement aux familles $\sL$ de classes de faisceaux inversibles sur les fibres de $X/S$.
\item Une famille $\Lambda$ de $\underline{\operatorname{Pic}}_{X/S}$ est quasi-compacte si et seulement si la famille $\sL$ correspondante est limitée.
\end{enumerate}
\end{prop}

\noindent\textit{En effet,} (i) résulte de (3.1) ; (ii), du fait suivant : si $s$ est un point de $S$ et $K$ est une extension algébriquement close de $k(s)$, alors les $K$-points de $\underline{\operatorname{Pic}}_{X/S}$ correspondent biunivoquement aux classes d'isomorphisme de faisceaux inversibles sur $X_K$.

Quant à (iii), supposons que la partie $\Lambda$ de $\underline{\operatorname{Pic}}_{X/S}$ soit quasi-compacte. Alors $\Lambda$ est contenue dans un ouvert $U$ de type fini, et l'inclusion $U \to \underline{\operatorname{Pic}}_{X/S}$ correspond à un faisceau inversible $L$ sur un $X_T$, où $T \to U$ est un morphisme surjectif (plat) de type fini. Enfin $T \to S$ est de type fini, donc $L$ limite $\sL$.

Réciproquement supposons que $\sL$ soit limitée par un faisceau cohérent $L$ sur un $X_T$, avec $T$ de type fini sur $S$. On peut supposer $L$ inversible (1.13 (i)) afin que $L$ définisse un morphisme $f : T \to \underline{\operatorname{Pic}}_{X/S}$. L'image de $f$ contient $\Lambda$ et elle est quasi-compacte, donc noethérienne par (i). Par suite $\Lambda$ est quasi-compact et on a (iii).

% ============================================================
% PAGE 636
% ============================================================
\begin{flushright}
$\fbox{\text{XIII}}$
\end{flushright}

Grâce à 3.2, la définition habituelle des morphismes de type fini entre deux schémas de Picard se généralise en la suivante :

\begin{defn}\label{def:3.3} Soient $X$ et $Y$ propres sur $S$ noethérien. Un morphisme entre les foncteurs de Picard
\[
\underline{\operatorname{Pic}}_{Y/S} \to \underline{\operatorname{Pic}}_{X/S}
\]
est dit \emph{de type fini} si toute famille limitée de classes de faisceaux inversibles sur les fibres de $X/S$ a comme image inverse une famille limitée de classes de faisceaux sur les fibres de $Y/S$.
\end{defn}

\begin{lem}\label{lem:3.4} Soient $X$, $Y$, $X'$, $Y'$ propres sur $S$ noethérien et $T \to S$ surjectif de type fini. Alors :
\begin{enumerate}[label={\roman*)}]
\item Pour qu'un morphisme $\Phi : \underline{\operatorname{Pic}}_{Y/S} \to \underline{\operatorname{Pic}}_{X/S}$ soit de type fini, il faut et il suffit que $\Phi_T : \underline{\operatorname{Pic}}_{Y/T} \to \underline{\operatorname{Pic}}_{X/T}$ soit de type fini.
\item Dans un diagramme commutatif
\[
\begin{tikzcd}
\underline{\operatorname{Pic}}_{Y'/S} \arrow[r] \arrow[d] & \underline{\operatorname{Pic}}_{Y/S} \arrow[d, "\Phi"] \\
\underline{\operatorname{Pic}}_{X'/S} \arrow[r] & \underline{\operatorname{Pic}}_{X/S}
\end{tikzcd}
\]
avec $\Phi'$ et $\Phi''$ de type fini, on a $\Phi$ de type fini.
\end{enumerate}
\end{lem}

\noindent\textit{En effet,} supposons $\Phi_T$ de type fini. Soit $\sL$ une famille sur les fibres de $X/S$, limitée par un $L$ sur un $X_P$. La famille $\sL_T = \sL \times_S T$ sur les fibres de $X_T/T$, définie de façon évidente, est limitée par l'image inverse de $L$ sur $X_{P \times_S T}$. Si $\Phi_T(\sL_T)$ est limité par un $M$ sur un $Y_Q$, alors $M$ limite aussi $\Phi(\sL)$.

La réciproque de (i) et l'assertion (ii) se démontrent de façon analogue.

\begin{thm}\label{thm:3.5} Soient $S$ un schéma noethérien, $X$ et $Y$ deux $S$-schémas propres et $f : X \to Y$ un morphisme. Si $f$ est surjectif, alors le morphisme induit entre les foncteurs de Picard
\[
f^* : \underline{\operatorname{Pic}}_{Y/S} \to \underline{\operatorname{Pic}}_{X/S}
\]
est de type fini.
\end{thm}

% ============================================================
% PAGE 638
% ============================================================
\begin{flushright}
$\fbox{\text{XIII}}$
\end{flushright}

\noindent\textit{En effet,} on peut supposer évidemment $S$ réduit et irréductible d'après 3.4 (i). Alors il existe un ouvert non-vide $U$ de $S$ tel que $\underline{\operatorname{Pic}}_{Y/S}$ et $\underline{\operatorname{Pic}}_{X/S}$ soient représentables et que la restriction de $f^*$ soit quasi-affine (XII 1.1), donc de type fini par 3.2 (i). Munissant $T = S - U$ de la structure réduite induite, on peut supposer que le morphisme $\underline{\operatorname{Pic}}_{Y/T} \to \underline{\operatorname{Pic}}_{X/T}$ induit par $f^*$ est de type fini par récurrence noethérienne ; d'où l'assertion, grâce à 3.4 (i) appliqué à $T \inj S$.

\begin{thm}\label{thm:3.6} Soient $S$ un schéma noethérien, $X$ un $S$-schéma propre et $\underline{\operatorname{Pic}}_{X/S}$ le foncteur de Picard. Alors pour tout entier $n \neq 0$, l'homomorphisme de puissance $n$-ième
\[
n : \underline{\operatorname{Pic}}_{X/S} \to \underline{\operatorname{Pic}}_{X/S}
\]
est un morphisme de type fini.
\end{thm}

\noindent\textit{En effet,} soit $f : X' \to X$ une présentation de Chow : $X'$ est projectif sur $S$ et $f$ est surjectif. Si l'on admet le théorème pour $X'$, alors en vertu de 3.5 et 3.4 (ii), le théorème pour $X$ en résulte. On peut donc supposer $X/S$ projectif.

Il existe un ouvert non-vide $U$ de $S$, un schéma intègre $S'$ et un morphisme fini et surjectif $S' \to U$, tels que toute composante irréductible $X'_j$ de $X_{S'}$, munie de la structure réduite induite, soit intègre sur $S'$ (XII 1.3). Munissant $T = S - U$ de la structure réduite induite, on peut admettre le théorème pour $X_T/T$ par récurrence noethérienne. Si l'on l'admet pour tous les $X'_j/S'$, il en résulte pour $\coprod_j X'_j/S'$ par le lemme facile suivant, donc pour $X_{S'}/S'$ en vertu de 3.5 et 3.4 (ii), enfin pour $X/S$ grâce à 3.4 (i).

\begin{lem}\label{lem:3.7} Soient $X = \coprod_i X_i$ de type fini sur $S$ noethérien et $\sF$ une famille de classes de faisceaux cohérents sur les fibres de $X/S$. Pour que $\sF$ soit limitée, il faut et il suffit que pour tout $i$ la famille $\sF_i$ induite sur les fibres de $X_i/S$ soit limitée.
\end{lem}

\noindent\textit{Dans la démonstration de 3.6,} on s'est ramené au cas où $X$ est projectif et intègre sur $S$ ; munissons $X/S$ d'un faisceau très ample $\cO_X(1)$. Soient $\sL$ une famille limitée de classes de faisceaux inversibles sur les fibres de $X/S$ et $L$ un faisceau inversible sur un $X_T$ avec $T$ de type fini sur $S$, qui limite $\sL$. Quitte à couper $T$ en morceaux, on peut supposer $L$ plat sur $T$, et quitte à remplacer $T$ par chacune de ses composantes connexes, on peut supposer $T$ connexe. Alors le polynôme en $q$ et $m$, $P(q,m) = \chi(L_t^{\otimes q}(m))$ ne dépend pas de $t \in T$ (EGA III 7.9.4).

% ============================================================
% PAGE 639
% ============================================================
\begin{flushright}
$\fbox{\text{XIII}}$
\end{flushright}

\noindent
Soit $\sL_n = \{M_t : \sL_t \text{ est la famille des classes des faisceaux inversibles } M_t \text{ tels que la classe de } M_t^{\otimes n} \text{ est dans } \sL\}$. Soit $\chi(M_t^{\otimes q}(m)) = \sum_i a_i(q) \binom{m+i}{i}$, où $a_i(q)$ est un polynôme en $q$ de degré $\leq r-i$. Pour $q, m$ variables, on a $\chi(M_t^{\otimes q}(m)) = P(q,m)$ et chaque $a_i(q)$ ne dépend donc pas du $M_t$ choisi ; par suite chaque polynôme $a_i(q)$ ne dépend pas non plus de $M_t$. Par conséquent la famille $\sL_n$ a un seul polynôme de Hilbert $\chi(M_t(m))$, donc est limitée en vertu de 2.11.

\begin{thm}\label{thm:3.8} Soient $S$ un schéma noethérien, $Y$ un $S$-schéma projectif muni d'un faisceau inversible relativement ample $\cO_Y(1)$, $X$ le schéma des zéros d'une section $\sigma$ de $\cO_Y(1)$, et $f : X \inj Y$ le morphisme d'inclusion. On suppose que pour tout point $s \in S$, et toute composante irréductible $Y_s^j$ de $Y_s$, les composantes irréductibles de $X_s \cap Y_s^j = V(\sigma|Y_s^j)$ soient de dimension $\geq 2$, (ce qui est le cas si les composantes irréductibles des fibres de $Y$ sont de dimension $\geq 3$). Alors le morphisme induit entre les foncteurs de Picard
\[
f^* : \underline{\operatorname{Pic}}_{Y/S} \to \underline{\operatorname{Pic}}_{X/S}
\]
est de type fini.

De plus, si $S$ est intègre, il existe un ouvert non-vide $U$ de $S$ tel que les foncteurs de Picard restreints à $U$ soient représentables et le morphisme induit
\[
\underline{\operatorname{Pic}}_{Y/S}|_U \to \underline{\operatorname{Pic}}_{X/S}|_U
\]
soit quasi-affine (de type fini).
\end{thm}

\noindent\textit{Démontrons d'abord la première assertion.} Comme dans la preuve de 3.6, on voit par récurrence noethérienne qu'on peut remplacer $S$ par un ouvert non-vide quelconque, et on se ramène facilement au cas où $S$ est intègre et $Y$ est intègre sur $S$, toujours en remplaçant $X$ par $X \times_Y Y'$ lorsqu'on remplace $Y$ par $Y'$ ; ceci fait, on se ramène de la même façon au cas où $Y$ est normal sur $S$, en utilisant à la place de XII 1.3 le lemme suivant.

\begin{lem}\label{lem:3.9} Soient $S$ un schéma noethérien intègre et $X$ un $S$-schéma, intègre sur $S$. Alors il existe un ouvert non-vide $U$ de $S$, un schéma intègre $S'$ fini et fidèlement plat sur $U$, un $S'$-schéma $X'$ normal et intègre sur $S'$, et un $S'$-morphisme fini surjectif $X' \to X_{S'}$.
\end{lem}

% ============================================================
% PAGE 640
% ============================================================
\begin{flushright}
$\fbox{\text{XIII}}$
\end{flushright}

\noindent\textit{En effet,} soient $\eta$ le point générique de $S$ et $K$ la clôture algébrique de $k(\eta)$. Il est clair que $\Spec(K)$ est la limite du système projectif filtrant $(S_{\alpha})$ des schémas affines intègres munis d'un morphisme fini et fidèlement plat $S_{\alpha} \to U_{\alpha}$, où $U_{\alpha}$ est un ouvert non-vide de $S$. Le normalisé de $X_{\eta}$ provient donc d'un $S_{\alpha}$-schéma $X'$ pour un $\alpha$, le morphisme fini et surjectif $X'_K \to X_{\eta}$ d'un $S_{\alpha}$-morphisme fini et surjectif $X' \to X_{S_{\alpha}}$ (EGA IV 8.8.2, 8.10.5). Enfin, comme $X'_K$ est normal et intègre, il existe un ouvert $S'$ de $S_{\alpha}$ tel que $X'_{S'}$ soit normal et intègre sur $S'$ (EGA IV 3.9.5).

\begin{rem*}\label{rem:3.10}
Conservant les notations et conditions de 3.9, supposons qu'il existe un morphisme résolvant $X'_K \to X_{\eta}$ (EGA IV 7.9.1), (par exemple si $\dim(X_{\eta}) \leq 2$, d'après Abhyankar). Alors un raisonnement similaire montre qu'il existe un morphisme fini et fidèlement plat $S' \to U$, un schéma $X'$ lisse sur $S'$ et un morphisme propre surjectif birationnel $X' \to X_{S'}$.
\end{rem*}

\noindent\textit{Dans la démonstration de la première partie de 3.8,} on s'est ramené au cas où $S$ est intègre et $Y$ est normal et intègre sur $S$, donc $Y$ est intègre. Si $\sigma = 0$, alors $X=Y$ et l'assertion est triviale. Sinon, quitte à restreindre $S$, on peut supposer que pour tout $s \in S$, on a $\sigma_s \neq 0$, donc $\sigma_s$ est $\cO_{Y_s}$-régulier et la dimension $r$ de $Y_s$ est $\geq 3$.

Soit alors $\sL$ une famille limitée de classes de faisceaux inversibles sur les fibres de $X/S$. Soit $M_s$ un faisceau inversible sur une fibre géométrique $Y_s$ de $Y/S$, avec le polynôme de Hilbert $\chi(M_s(n)) = \sum_i a_i \binom{n+i}{i}$ ; alors le polynôme de Hilbert de $M_s|_{X_s}$ est $\sum_{i=0}^{r-1} a_{i+1} \binom{n+i}{i}$. Lorsque la classe de $M_s$ parcourt $(f^*)^{-1}(\sL)$, c-à-d la classe de $f^*M_s$ est dans $\sL$, les trois coefficients $a_{r-1}$, $a_{r-2}$, $a_{r-3}$ restent bornés grâce à 1.13.

% ============================================================
% PAGE 641 -- 642
% ============================================================
\begin{flushright}
$\fbox{\text{XIII}}$
\end{flushright}

\noindent
Soit $H = \cO_Y(m)$ un multiple de $\cO_Y(1)$ qui est très ample. Alors dans les polynômes de Hilbert $\chi(M_s^{\otimes q}(n)) = \sum_i a_i^{(m)} \binom{n+i}{i}$ relativement à $H$ des $M_s$ dans $(f^*)^{-1}(\sL)$, les trois coefficients $a_{r-1}^{(m)}$, $a_{r-2}^{(m)}$, $a_{r-3}^{(m)}$ restent bornés grâce à 2.10. Par conséquent $\sL$ est limitée en vertu de 2.11 ; d'où la première partie de 3.8.

Quant à la deuxième partie de 3.8, la représentabilité ``générique'' des foncteurs de Picard a été établie dans XII 1.2, ainsi que leur séparation. Par un raisonnement similaire à celui de XII 2, utilisant 3.9 et XII 1.1, on se ramène au cas $Y$ normal sur $S$. Or, pour que le morphisme $\underline{\operatorname{Pic}}_{Y/S} \to \underline{\operatorname{Pic}}_{X/S}$ soit quasi-affine, il suffit qu'il soit (séparé et) quasi-fini (EGA IV 8.12.8), ce qui se vérifie sur les fibres géométriques, et résulte donc du lemme suivant :

\begin{lem}\label{lem:3.11} Soient $k$ un corps algébriquement clos, $Y$ un $k$-schéma normal, projectif muni d'un faisceau ample $\cO_Y(1)$, et $X$ le schéma des zéros d'une section $\sigma$ de $\cO_Y(1)$. On suppose que $Y$ soit de dimension $\geq 3$ (resp.\@ $\geq 2$). Alors le morphisme $\Pic_Y \to \Pic_X$ (resp.\@ le morphisme $\Pic^0_Y \to \Pic^0_X$ (voir 4.1)) est de type fini et a un noyau $N$ qui est un groupe fini (non-réduit en général).
\end{lem}

\noindent\textit{En effet,} on peut évidemment supposer $\sigma \neq 0$ ; donc $\sigma$ est régulière, $Y$ étant intègre. Le morphisme envisagé est de type fini grâce à la première partie de 3.8 (resp.\@ grâce à 4.7 (ii)). Par suite le noyau $N$ est un groupe algébrique, donc contenu dans $\Pic^0_Y$, donc est le même dans les deux cas envisagés. Comme $\Pic^0_Y$ est propre, $Y$ étant normal (EGA 236--2.13), et comme $N$ est un sous-groupe, nécessairement fermé (SGA 3 VI, 1.9.2), pour que $N$ soit fini, il suffit qu'il ne contienne pas le sous-groupe de type $\bmu_{\ell} = \bZ/\ell\bZ$ où $(\ell, \car(k)) = 1$.

% ============================================================
% PAGE 643
% ============================================================
\begin{flushright}
$\fbox{\text{XIII}}$
\end{flushright}

\noindent
Considérons la théorie de Kummer, c'est-à-dire, le diagramme commutatif suivant, provenant de la suite exacte
\[
0 \to \bmu_{\ell} \to \bG_m \xrightarrow{\;\ell\;} \bG_m \to 0
\]
de faisceaux étales sur $Y$ (resp.\@ sur $X$) :
\[
\begin{tikzcd}
\HH^1(Y, \bmu_{\ell}) \arrow[r] \arrow[d] & \HH^1(Y, \bG_m) \arrow[r, "\ell"] \arrow[d] & \HH^1(Y, \bG_m) \arrow[d] \\
\HH^1(X, \bmu_{\ell}) \arrow[r] & \HH^1(X, \bG_m) \arrow[r, "\ell"'] & \HH^1(X, \bG_m)
\end{tikzcd}
\]
Par le ``théorème 90'' de Hilbert, on a $\HH^1(Y, \bG_m) = \Pic_Y(k)$ (resp.\@ $\HH^1(X, \bG_m) = \Pic_X(k)$). Comme $X$ est propre et intègre, on a $\HH^0(X, \bG_m) = k^*$ ; donc $\HH^1(Y, \bmu_{\ell}) \to \HH^1(Y, \bG_m)$ est injectif.

Il suffit donc de voir que $\HH^1(Y, \bZ/\ell\bZ) \to \HH^1(X, \bZ/\ell\bZ)$ est injectif, ce qui résulte du fait que $\pi_1(X) \to \pi_1(Y)$ est surjectif, ou encore (SGA 1 V 6.9) que tout revêtement étale connexe $\tilde{Y}$ de $Y$ de groupe $\bZ/\ell\bZ$ a une restriction $\tilde{X} = X \times_Y \tilde{Y}$ qui est aussi connexe, ce qui est une conséquence bien connue du théorème de Bertini et du théorème de connexion de Zariski (SGA 1 X 2.10).

\begin{rem*}\label{rem:3.12}
Conservons les conditions de 3.11. Alors :
\begin{enumerate}[label={\roman*)}]
\item Le noyau $N$ est un groupe fini unipotent (SGA 3 XVII 1.1), puisqu'il ne contient pas le sous-groupe de type multiplicatif $\bG_m$ pour aucun $n$ (SGA 3 XVII 4.6.1 vi), ce qu'on prouve de même, utilisant au lieu de la théorie de Kummer le fait suivant, dû à Grothendieck (SGA 1 XI, à la dernière page) : si $X$ est un schéma propre sur un corps $k$, avec $\HH^0(X, \cO_X) = k$, alors il existe un isomorphisme de $k$-schémas en groupes $\Hom_k(\bZ/n\bZ, \Pic^0_{X/k}) \simeq \HH^1(X, \bZ/n\bZ)$ ; par conséquent, si $k$ est algébriquement clos, on a $\Hom_k(\bZ/n\bZ, \Pic^0_{Y/k}) \simeq \HH^1(Y, \bZ/n\bZ)$.
\item Si $k$ est de caractéristique nulle, alors le noyau $N$ est nul, puisque tout groupe fini unipotent est alors nul (SGA 3 XVII 1.7). On retrouve ainsi à peu près le ``Vanishing theorem'' de Kodaira sous la forme donnée par Mumford [2].
\end{enumerate}
\end{rem*}

% ============================================================
% PAGE 644
% ============================================================
\begin{flushright}
$\fbox{\text{XIII}}$
\end{flushright}

\begin{thm}\label{thm:3.13} Soient $S$ un schéma noethérien et $X$ un $S$-schéma projectif, muni d'un faisceau inversible relativement ample $H = \cO_X(1)$. On suppose que les fibres géométriques de $X/S$ soient intègres et de dimension $r$. Alors pour une famille $\sL$ de classes de faisceaux inversibles sur les fibres de $X/S$, les conditions suivantes sont équivalentes :
\begin{enumerate}[label={\roman*)}]
\item $\sL$ est une famille limitée.
\item Dans les polynômes de Hilbert $\chi(L_s(n)) = \sum_{i=0}^{r} a_i(s) \binom{n+i}{i}$ des $L_s \in \sL$, le coefficient $a_{r-1}$ reste borné et le coefficient $a_{r-2}$ reste minoré.
\item Lorsque $L_s$ parcourt $\sL$, le degré $d(L_s) = \langle c_1(L_s) \cdot c_1(H_s)^{r-1} \rangle$ reste borné et le nombre d'intersection $\langle c_1(L_s)^2 \cdot c_1(H_s)^{r-2} \rangle$ reste minoré.
\item Il existe un entier $N$ tel que l'on ait
\[
\sL^{\otimes N} < \sL < \sL^{\otimes -N}
\]
au sens de la relation d'ordre définie par la relation d'amplitude ; c'est-à-dire, pour tout $L_s \in \sL$, $L_s^{\otimes N} \otimes H_s^{\otimes -1}$ et $L_s^{\otimes -N} \otimes H_s$ sont amples.
\end{enumerate}
\end{thm}

% ============================================================
% PAGE 645
% ============================================================
\begin{flushright}
$\fbox{\text{XIII}}$
\end{flushright}

\noindent\textit{Démontrons que (i) implique (ii), (iii), et (iv).} Supposons donc que $\sL$ soit limitée par un faisceau inversible sur un $X_T$. Alors $\sL$ satisfait à (iv) par EGA II 4.6.12. Supposons $L$ plat sur $T$. Alors le polynôme en $m$ et $n$, $\chi(L_t^{\otimes m}(n))$ ne dépend que de la composante connexe de $T$ (EGA III 7.9.4). Par suite $\sL$ satisfait à (ii), et compte tenu de 2.1, à (iii).

Notons ensuite que la famille des classes des $H_K$ (resp.\@ $\cO_{X_K}$) est limitée, donc satisfait à (iii) (resp.\@ à (ii)).

Supposons que $\sL$ satisfasse à (iv). Alors
\[
\langle N \cdot c_1(H_K) + c_1(L_K) \rangle \cdot c_1(H_K)^{r-2} > 0, \text{ d'où } d(L_K) > -N d(H_K),
\]
et de même, $N d(H_K) > d(L_K)$. De plus, $\langle N \cdot c_1(H_K) + c_1(L_K) \rangle^2 \cdot c_1(H_K)^{r-2} > 0$, d'où $\langle c_1(L_K)^2 \cdot c_1(H_K)^{r-2} \rangle > -2N d(L_K) - N^2 d(H_K)$. Donc $\sL$ satisfait à (iii).

\noindent\textit{Démontrons enfin que (ii) (resp.\@ (iii)) implique (i).} Grâce à 2.10 et 2.3, on peut remplacer $H$ par un multiple ; supposons donc $H$ très ample. Évidemment on peut couper $S$ en morceaux ; supposons donc qu'il existe une section linéaire $Y$ de $X$, qui est intègre sur $S$, à fibres de dimension $2$. En vertu de 3.8, on peut remplacer $X$, $\sL$ par $Y$, $\sL|Y$, (compte tenu de 1.7 (resp.\@ la formule de projection), qui implique que la condition (ii) (resp.\@ (iii)) sur $X$, $\sL$ est équivalente à (ii) (resp.\@ à (iii)) sur $Y$, $\sL|Y$) ; supposons donc $r=2$. Si $\delta = \sup\{d(L_K)\}$ est $> 0$, on peut remplacer $\sL$ par la famille des classes des $L_K(-\delta)$ grâce à 2.3 et 2.10 ; supposons donc $\delta \leq 0$.

Si dans les polynômes de Hilbert $\sum_{i=0}^{2} a_i \binom{n+i}{i}$, $a_1$ reste borné et $a_0$ reste minoré, alors les trois coefficients restent bornés grâce à 2.3 et 2.7 ; donc $\sL$ est limitée en vertu de 2.11 ; ainsi (ii) implique (i). D'autre part, (iii) implique (ii) grâce à 2.3 et 2.8.

% ============================================================
% PAGE 646
% ============================================================
\section*{4. Théorèmes de finitude pour $\Pic^{\tau}_{X/S}$}
\addcontentsline{toc}{section}{4. Théorèmes de finitude pour $\Pic^{\tau}_{X/S}$}

\begin{flushright}
$\fbox{\text{XIII}}$
\end{flushright}

\smallskip
\noindent\textbf{4.1.} Soient $S$ un schéma et $G$ un foncteur en groupes commutatifs sur $S$. On désigne par $G^0$ (resp.\@ $G^{\tau}$) le sous-foncteur en groupes dont les $T$-points $f : T \to G$ sont ceux de $G$ tels que pour tout corps algébriquement clos $K$ et tout $K$-point $t$ de $T$, il existe un $k$-schéma connexe $Z$, deux $K$-points $z$, $z_0$ de $Z$ et un morphisme $g : Z \to G$ tels que $g(z) = f(t)$ et $g(z_0) = 0$ dans $G(K)$ (resp.\@ il existe un entier non-nul $n$ tel que $n \cdot f(t) \in G^0(K)$).

Il est évident que la formation de $G^0$ (resp.\@ $G^{\tau}$) commute à tout changement de base, qu'un morphisme de foncteurs en groupes $G \to G'$ induit un morphisme $G^0 \to G'^0$ (resp.\@ $G^{\tau} \to G'^{\tau}$), et que si $S$ est un corps et $G$ est représenté par un schéma en groupes localement de type fini, alors $G^0$ est représenté par la composante connexe de l'élément neutre, qui est un sous-schéma ouvert (et fermé), de type fini, et géométriquement irréductible (SGA 3 VI, 2.4).

\begin{lem}\label{lem:4.2} Soit $\Phi : G \to G'$ un morphisme entre deux foncteurs en groupes commutatifs sur un schéma $S$. On suppose que pour tout corps algébriquement clos $K$ et tout morphisme $\Spec(K) \to S$, $G_K$ (resp.\@ $G'_K$) soit représentable par un schéma en groupes localement de type fini, et $\Phi_K$ soit représentable par un morphisme de type fini, (c-à-d, quasi-compact). Alors $\Phi(G^{\tau}) \subset G'^{\tau}$.
\end{lem}

% ============================================================
% PAGE 647
% ============================================================
\begin{flushright}
$\fbox{\text{XIII}}$
\end{flushright}

\noindent\textit{En effet,} on a déjà noté l'inclusion $\Phi(G^0) \subset G'^0$. Pour l'inclusion en sens inverse, on note que $g \in G(K)$ est dans $G^{\tau}(K)$ si et seulement si la suite des $ng$ ($n \in \bN$) est quasi-compacte ; comme $Z$ est quasi-compact, il suffit pour ceci que la suite des $\Phi(ng) = n\Phi(g)$ ait la même propriété ; d'où l'assertion.

\begin{lem}\label{lem:4.3} Soient $X$, $Y$ propres sur un schéma $S$ et $f : X \to Y$ un morphisme surjectif (resp.\@ le morphisme d'inclusion du schéma des zéros d'une section d'un faisceau ample $\cO_Y(1)$). Alors
\[
(f^*)^{-1}(\Pic^{\tau}_{X/S}) = \Pic^{\tau}_{Y/S}.
\]
\end{lem}

\noindent\textit{En effet,} l'assertion est conséquence formelle de 4.2, (XII 1.2), 3.1, et 3.5 (resp.\@ 3.8).

\smallskip
\noindent\textbf{4.4.} Soient $k$ un corps algébriquement clos, $X$ un $k$-schéma propre, et $L$ un faisceau inversible sur $X$. On dit que $L$ est \emph{algébriquement équivalent à} $0$ (resp.\@ $\tau$-\emph{équivalent à} $0$) s'il existe un $k$-schéma connexe de type fini $T$, un faisceau inversible $M$ sur $X_T = X \times_k T$, et deux $k$-points $t$, $t_0$ de $T$ tels que $L = M_t$ et $\cO_X \simeq M_{t_0}$ (resp.\@ s'il existe un entier non nul $n$ tel que $L^{\otimes n}$ soit algébriquement équivalent à $0$), ou ce qui revient facilement au même, si le $k$-point $\lambda$ de $\Pic_{X/k}$ qui correspond à $L$, est dans $\Pic^0_{X/k}(k)$ (resp.\@ dans $\Pic^{\tau}_{X/k}(k)$).

Rappelons aussi que $L$ est dit \emph{numériquement équivalent à} $0$ si pour toute courbe intègre fermée $Y$ dans $X$, le nombre d'intersection $\langle c_1(L) \cdot Y \rangle$ est $0$ (voir 2.1 ou X 4.5).

% ============================================================
% PAGE 648
% ============================================================
\begin{flushright}
$\fbox{\text{XIII}}$
\end{flushright}

\begin{prop}\label{prop:4.5} \textup{(Comportement fonctoriel).} Soit $f : X' \to X$ un morphisme entre deux $k$-schémas propres, et soit $L$ un module inversible sur $X$. Alors :
\begin{enumerate}[label={\roman*)}]
\item Si $L$ est $\tau$-équivalent à $0$ (resp.\@ numériquement équivalent à $0$), alors $f^*L$ l'est.
\item Réciproquement, lorsque $f$ est surjectif, si $f^*L$ est $\tau$-équivalent à $0$ (resp.\@ numériquement équivalent à $0$), alors $L$ l'est.
\end{enumerate}
\end{prop}

\noindent\textit{En effet,} l'assertion (i) (resp.\@ (ii)) pour la $\tau$-équivalence résulte aussitôt de 4.1 (resp.\@ 4.3). D'autre part, soient $Y'$ une courbe intègre fermée dans $X'$, et $Y = f(Y')$ son image, munie de la structure réduite induite. Lorsque $\dim(Y) = 1$, on a la formule de projection (X 4.4.4)
\[
\text{(4.5.1)} \qquad \langle c_1(f^*L) \cdot Y' \rangle = [k(Y') : k(Y)] \, \langle c_1(L) \cdot Y \rangle
\]
où $k(Y')$ (resp.\@ $k(Y)$) est le corps de fonctions de $Y'$ (resp.\@ $Y$), et lorsque $\dim(Y) = 0$, on a $(f^*L)|_{Y'} \simeq \cO_{Y'}$. Par conséquent, si $L$ est numériquement équivalent à $0$, alors $\langle c_1(f^*L) \cdot Y' \rangle = 0$ ; d'où (i).

Réciproquement, lorsque $f$ est surjective, étant donnée une courbe intègre fermée $Y$ dans $X$, on peut trouver une courbe intègre fermée $Y'$ dans $X'$ telle que $Y = f(Y')$ : soient $y$ le point générique de $Y$, $y'$ un point de la fibre $f^{-1}(y)$ algébrique sur $k(y)$, et $Y'$ l'adhérence de $y'$. Alors 4.5.1 montre que $\langle c_1(f^*L) \cdot Y' \rangle = 0$ implique $\langle c_1(L) \cdot Y \rangle = 0$ ; d'où (ii).

\begin{thm}\label{thm:4.6} \textup{(La caractérisation numérique de $\Pic^{\tau}_{X/k}$).} Soient $k$ un corps algébriquement clos, $X$ un $k$-schéma propre, et $L$ un faisceau inversible sur $X$. Alors les conditions suivantes sont équivalentes :
\begin{enumerate}[label={\roman*)}]
\item $L$ est $\tau$-équivalent à $0$.
\item Pour tout faisceau cohérent $F$ sur $X$, on a $\chi(F \otimes L) = \chi(F)$.
\item $L$ est numériquement équivalent à $0$.
\end{enumerate}
Si $X/k$ est projectif et muni d'un module inversible ample $H = \cO_X(1)$, les conditions précédentes sont aussi équivalentes aux suivantes :
\begin{enumerate}[label={\roman*)},resume]
\item Pour tout entier $m$ (positif ou négatif), $L^{\otimes m}$ est ample.
\item (Si $X$ est irréductible). Pour tout couple d'entiers $m$, $n$, on a $\chi(L^{\otimes m}(n)) = \chi(\cO_X(n))$.
\item (Si $X$ est irréductible de dimension $r \geq 2$). On a $\langle c_1(L) \cdot H^{r-1} \rangle = 0$, $\langle c_1(L)^2 \cdot H^{r-2} \rangle = 0$.
\end{enumerate}
\end{thm}

% ============================================================
% PAGE 649
% ============================================================
\begin{flushright}
$\fbox{\text{XIII}}$
\end{flushright}

\noindent\textit{Démontrons que (i) implique (ii).} Appliquant le lemme de dévissage (EGA III 3.1) et 4.5 (i), on se ramène au cas où $X$ est intègre et $F = \cO_X$. Alors on doit prouver que le polynôme $\chi(L^{\otimes m})$ est constant. Pour ceci, on peut remplacer $L$ par un multiple, donc supposer $L$ algébriquement équivalent à $0$. Dans ce cas, si $L$ est déformé en $\cO_X$ par la famille algébrique définie par un faisceau inversible $M$ sur un $X \times T$ où $T$ est connexe, de type fini, le polynôme $\chi(M_t)$ ne dépend pas de $t \in T$ (EGA III 7.9.5) ; d'où l'implication.

Soit $Y$ une courbe intègre fermée dans $X$. Faisant $F = \cO_Y$ dans (ii), on trouve $\chi(\cO_Y) = \chi(L|_Y)$. Donc la formule (``théorème de Riemann'') $\chi(L|_Y) = \langle c_1(L) \cdot Y \rangle + \chi(\cO_Y)$ donne $\langle c_1(L) \cdot Y \rangle = 0$ ; d'où (ii) implique (iii).

Soit $f : X' \to X$ une présentation de Chow : $X'$ est projectif et $f$ est surjectif. Si l'on admet que (iii) implique (i) pour $X'$, alors

% ============================================================
% PAGE 650
% ============================================================
\begin{flushright}
$\fbox{\text{XIII}}$
\end{flushright}

\noindent
en vertu de 4.5, l'implication pour $X$ en résulte. Supposons donc dorénavant $X$ projectif, muni d'un faisceau ample $H$.

Soient $X_i$ les composantes irréductibles de $X$, munies de la structure réduite induite. Appliquant 4.5 et EGA III 2.6.2 au morphisme fini et surjectif $\coprod_i X_i \to X$, on ramène la démonstration que (iii) implique (i) et (iv) au cas où $X$ est intègre.

Si $X$ est irréductible, il est évident que (ii) implique (v), que (iii) implique (vi), et compte tenu de 2.1, que (v) implique (vi), et que (vi) pour $X$ est équivalent à (vi) pour $X_{\red}$ (X 4.4.1).

Il reste donc à voir que (vi) implique (i) et (iv), et que (iv) implique (i), en supposant $X$ intègre. Mais alors, (vi) (resp.\@ (iv)) implique que la famille $\sL$ des classes des $L^{\otimes m}$ ($m \in \bN$) satisfait à 3.13 (iii) (resp.\@ (iv)), donc est limitée. Il existe donc un faisceau inversible $L$ sur un $X_T$, avec $T$ de type fini sur $k$, qui limite $\sL$. Soient $p$, $q$ deux entiers différents tels que $L^{\otimes p}$ et $L^{\otimes q}$ correspondent à la même composante connexe de $T$ ; alors $L^{\otimes (p-q)}$ est donc algébriquement équivalent à $0$ ; d'où (vi) (resp.\@ (iv)) implique (i).

Enfin, on a de même que (vi) implique qu'il existe un entier $N$ tel que pour tout entier $m$, $L^{\otimes m} \otimes H^{\otimes N}$ soit ample. Alors pour tout $m$, $(L^{\otimes m} \otimes H)^{\otimes N}$ est ample ; d'où que (vi) implique (iv), ce qui achève la démonstration.

\begin{thm}\label{thm:4.7} Soient $S$ un schéma noethérien, et $X$ un $S$-schéma propre. Alors :
\begin{enumerate}[label={\roman*)}]
\item Le morphisme d'inclusion $\Pic^{\tau}_{X/S} \to \Pic^{\mathrm{num}}_{X/S}$ est représentable par une immersion ouverte.
\item Si $X$ est projectif et intègre sur $S$, alors $\Pic^{\tau}_{X/S} \to \underline{\operatorname{Pic}}_{X/S}$
\end{enumerate}
\end{thm}

\begin{center}
\rule{0.5\textwidth}{0.4pt}
\end{center}


\clearpage
% ===== Source block: SGA 6 pages 651 702 =====
\phantomsection
\addcontentsline{toc}{section}{SGA 6 pages 651 702}
\section*{SGA 6 pages 651 702}
\setcounter{page}{651}

\thispagestyle{fancy}

\noindent\rule{\textwidth}{0.5pt}


% ============================================================
\section*{Expos\'e XIII (suite et fin)}
\addcontentsline{toc}{section}{Expos\'e XIII --- Quotient de Pic par sous-sch\'emas de Pic (suite et fin)}

\setcounter{section}{4}
\setcounter{subsection}{0}

\subsection*{4.7. Compl\'ements}
\addcontentsline{toc}{subsection}{4.7. Compl\'ements}

\begin{theoreme}[4.7]
Soit $f : X \to S$ propre et plat, avec $S$ noeth\'erien. Alors :

\emph{(i)} Le sous-foncteur $\Pic^{\tau}_{X/S}$ de $\Pic_{X/S}$ des classes de faisceaux inversibles num\'eriquement \'equivalents \`a $0$ sur les fibres est repr\'esentable par un sous-sch\'ema ouvert de $\Pic_{X/S}$.

\emph{(ii)} Supposons de plus que $f$ soit de Cohen-Macaulay, et que les fibres g\'eom\'etriques de $f$ soient int\`egres. Alors le sous-foncteur $\Pic^{\num}_{X/S}$ de $\Pic_{X/S}$ des classes de faisceaux inversibles alg\'ebriquement \'equivalents \`a $0$ sur les fibres est aussi repr\'esentable par un sous-sch\'ema ferm\'e.

\emph{(iii)} $\Pic^{\tau}_{X/S}$ est de type fini sur $S$, c'est-\`a-dire, la famille des classes des faisceaux inversibles $\tau$-\'equivalents \`a $0$ sur les fibres de $X/S$ est limit\'ee.
\end{theoreme}

En effet, dans (i) (resp.\ (ii)), on doit traiter le cas d'un morphisme $T \to \Pic_{X/S}$, o\`u $T$ est un $S$-sch\'ema quelconque, et montrer que l'image inverse de $\Pic^{\tau}_{X/S}$ (resp.\ $\Pic^{\num}_{X/S}$) est repr\'esentable par un sous-sch\'ema ouvert (resp.\ et ferm\'e) de $T$. Il suffit de prendre $T = \Spec(B)$ sur un ouvert affine $\Spec(A)$ de $S$. Alors $T$ est la limite du syst\`eme projectif filtrant des $S$-sch\'emas affines $T_i = \Spec(B_i)$, o\`u $B_i$ parcourt les sous-$A$-alg\`ebres de type fini de $B$. En vertu de 3.1, le morphisme en question provient d'un morphisme $T_i \to \Pic_{X/S}$. On peut donc supposer $T$ de type fini sur $S$.

Le morphisme $T \to \Pic_{X/S}$ correspond \`a un morphisme f.p.p.f.\ $R \to T$ et un faisceau inversible $L$ sur $X_R$. Par descente f.p.p.f., on peut remplacer $T$ par $R$. Apr\`es le changement de base $R \to S$, compte tenu de la compatibilit\'e de la formation de $\Pic^{\tau}_{X/S}$ avec tout changement de base (4.1) on est donc ramen\'e dans (i) et (ii) \`a \'etudier le cas d'une section $S \to \Pic_{X/S}$, d\'efinie par un faisceau inversible $L$ sur $X$.

On doit prouver que l'ensemble $S_0$ des $s \in S$ tels que $L_s$ soit $\tau$-\'equivalent \`a $0$ est ouvert, et aussi ferm\'e sous les conditions de (ii); comme $S$ est noeth\'erien, ceci veut dire trois choses : pour un $s \in S_0$,

\begin{enumerate}[label=(\alph*)]
\item toute g\'en\'erisation de $s$ reste dans $S_0$;
\item toute sp\'ecialisation de $s$, voisine de $s$, reste dans $S_0$;
\item sous les conditions de (ii), toute sp\'ecialisation de $s$ reste dans $S_0$.
\end{enumerate}

\medskip
\noindent\textbf{Preuve de (a), (b), (c) et (iii).}

Soit $f : X' \to X$ une pr\'esentation de Chow : $X'$ est projectif sur $S$; et $f$ est surjectif. Si l'on admet (a), (b), (c) et (iii) pour $X'$, alors ils en r\'esultent pour $X$ en vertu de 4.3 et 3.5. On peut donc supposer $X$ projectif sur $S$.

Soit $H$ ample pour $X/S$. Si pour un $s \in S$, $(L \otimes H)_s$ est ample, alors $(L \otimes H)_t$ est ample pour toute g\'en\'erisation $t$ de $s$ (EGA III 4.7.1). Gr\^ace \`a l'\'equivalence de 4.6 (i) et (iv), on a donc \'etabli (a).

Pour (b), (c) et 4.7 (iii), on peut supposer $S$ int\`egre avec $s$ point g\'en\'erique dans (b) et (c) (Pour (b) et (c), c'est clair, et pour (iii), cela r\'esulte de 3.4 (i) appliqu\'e au morphisme $\coprod S_i \to S$, o\`u les $S_i$ sont les composantes irr\'eductibles de $S$, munies de la structure r\'eduite induite).

Il existe un sch\'ema int\`egre $S'$ et un morphisme f.p.p.f.\ $S' \to S$, tels que toute composante irr\'eductible $X_j$ de $X_{S'}$, munie de la structure r\'eduite induite soit int\`egre sur $S'$ (XII 1.3). Pour prouver (b) et (iii), il suffit de le faire pour $X_{S'}/S'$, (dans (b), par descente; dans (iii) par r\'ecurrence noeth\'erienne et 3.4 (i)).

Enfin, consid\'erons le morphisme surjectif canonique $\coprod X'_j \to X_{S'}$. Alors (b) et (iii) pour $X_{S'}/S'$ r\'esultent de (c) et (iii) pour les $X'_j/S'$, en vertu de 4.3, 3.5 et 3.7.

On s'est donc ramen\'e \`a prouver (c) et (iii), en supposant $S$ int\`egre de point g\'en\'erique $s$, et $X$ projectif et int\`egre sur $S$. Soit $\cO_X(1)$ un faisceau tr\`es ample pour $X/S$. Alors pour tout $m$, $n$, $\chi(L^m(n))$ et $\chi(\cO_X(n))$ ne d\'ependent pas de $t \in S$ (EGA III 7.9.4). Donc (c) r\'esulte de 4.6 (i) $\Leftrightarrow$ (v), qui implique aussi que la famille des classes des faisceaux inversibles $\tau$-\'equivalents \`a $0$ sur les fibres de $X/S$ n'a qu'un seul polyn\^ome de Hilbert; donc (iii) r\'esulte de 2.11.

\begin{remarque}[4.8]
Soit $X$ propre et int\`egre sur $S$ noeth\'erien. Sans hypoth\`ese de projectivit\'e sur $X/S$, on ignore s'il reste vrai que $\Pic^{\tau}_{X/S} \hookrightarrow \Pic_{X/S}$ est une immersion ferm\'ee, sauf lorsque $X$ est normal sur $S$ et $\Pic_{X/S}$ est repr\'esentable (FGA 236, 2.3).
\end{remarque}

\subsection*{5. Th\'eor\`emes de finitude du type ``N\'eron-Severi''}
\addcontentsline{toc}{subsection}{5. Th\'eor\`emes de finitude du type ``N\'eron-Severi''}

\begin{theoreme}[5.1]
Soit $X$ propre sur $S$ noeth\'erien; alors les groupes de N\'eron-Severi $\NS(X_{\ov{K}}) = \Pic_{X_{\ov{K}}/\ov{K}}(\ov{K}) / \Pic^{\tau}_{X_{\ov{K}}/\ov{K}}(\ov{K})$ des fibres g\'eom\'etriques $X_{\ov{K}}$ de $X/S$ sont de type fini, et leur rang et l'ordre de leur sous-groupe de torsion $\Tors(\NS(X_{\ov{K}}))$ restent major\'es.
\end{theoreme}

En effet, par r\'ecurrence noeth\'erienne, on peut se borner \`a prouver 5.1 en y rempla\c{c}ant $S$ par un ouvert non-vide arbitraire de $S_{\mathrm{red}}$; ainsi on peut supposer $\Pic^{\tau}_{X/S}$ repr\'esentable par un $S$-sch\'ema de type fini (XII 1.2 et 4.7 (iii)). Alors l'ordre de $\Tors(\NS(X_{\ov{K}}))$ est \'egal au nombre de composantes connexes de $(\Pic^{\tau}_{X/S})_{\ov{K}}$, nombre qui reste le m\^eme au-dessus d'un ouvert non-vide de $S$ (EGA IV 9.7.9).

Il reste \`a majorer le rang de $\NS(X_{\ov{K}})$. Comme dans la d\'emonstration de 4.7 (iii), on se ram\`ene facilement au cas $S$ int\`egre, et $X$ projectif et int\`egre sur $S$. Soit $r$ la dimension commune des fibres de $X/S$. Si $r$ est $0$ (resp.\ $1$), alors ce rang est \'evidemment $0$ (resp.\ $1$). Si $r$ est $\geq 2$, on va se ramener au cas $r = 2$ et $X$ lisse sur $S$.

Quitte \`a restreindre $S$ on peut trouver une section lin\'eaire $Y$ de $X$ qui est encore int\`egre sur $S$, \`a fibres de dimension $2$; en vertu de 4.3, on peut remplacer $X$ par $Y$. Enfin, en vertu de 4.3 et 3.10 (utilisant la r\'esolution des singularit\'es des surfaces), quitte \`a remplacer $S$ par un rev\^etement plat d'un ouvert non-vide de $S$, on peut remplacer $X$ par un $S$-sch\'ema lisse.

Maintenant, gr\^ace \`a 4.6 (i) $\Rightarrow$ (iii) et au th\'eor\`eme de finitude et de sp\'ecialisation des groupes de cohomologie \'etale (SGA 4 XVI 2.2), 5.1 r\'esulte de la proposition suivante.

\begin{proposition}[5.2]
Soient $k$ un corps alg\'ebriquement clos et $X$ un $k$-sch\'ema propre et lisse. Alors le groupe $C^p(X)$ des cycles alg\'ebriques de codimension $p$ sur $X$, modulo l'\'equivalence num\'erique, est libre de rang $\leq b^{2p}$, le nombre de Betti.
\end{proposition}

En effet, si l'on oublie l'action des racines de l'unit\'e, la cohomologie \'etale $H^i_{\mathrm{et}}(X, \mathbb{Q}_\ell)$ $= \bigl[\varprojlim H^i_{\mathrm{et}}(X, \mathbb{Z}/\ell^n\mathbb{Z})\bigr] \otimes_{\mathbb{Z}_\ell} \mathbb{Q}_\ell$ (pour $\ell \neq \car(k)$) devient une cohomologie de Weil : elle satisfait \`a la dualit\'e de Poincar\'e et \`a la formule de K\"unneth; et elle re\c{c}oit un homomorphisme fonctoriel, non nul de l'anneau des cycles alg\'ebriques modulo l'\'equivalence rationnelle $C^p(X) \to H^{2p}_{\mathrm{et}}(X, \mathbb{Q}_\ell(p))$ (voir X 7.13).

Soient $a_1, \dots, a_m \in C^{r-p}(X)$, o\`u $r = \dim(X)$, des \'el\'ements qui s'en vont sur une base du $\mathbb{Q}_\ell$-espace vectoriel dans $H^{2(r-p)}_{\mathrm{et}}(X, \mathbb{Q}_\ell(r-p))$ engendr\'e par les cycles alg\'ebriques. Consid\'erons l'homomorphisme
\[
\alpha : C^p(X) \to \mathbb{Q}_\ell^m
\]
d\'efini par $\alpha(b) = (\langle b.a_1 \rangle, \dots, \langle b.a_m \rangle)$, o\`u $\langle \; \rangle$ d\'esigne le nombre d'intersection. Alors on voit facilement que l'image de $\alpha$ est $C^p(X)$.

\begin{remarque}[5.3]
Soient $k$ un corps alg\'ebriquement clos, et $X$ une surface propre et lisse sur $k$. Alors le deuxi\`eme nombre de Betti $b^2(X)$ peut \^etre interpr\'et\'e comme \'etant $c_2(X) - 2 + 2b^1(X)$, o\`u $b^1(X) = 2 \dim(\Pic^0_{X/k})$ et $c_2(X)$ est la deuxi\`eme classe de Chern $=$ la caract\'eristique d'Euler-Poincar\'e $=$ le nombre d'intersection de la diagonale de $X \times X$ avec elle-m\^eme; ainsi $b^2(X)$ a une interpr\'etation ind\'ependante de la cohomologie. L'in\'egalit\'e $\rang(C^1(X)) \leq b^2(X)$ dans cette forme pour $\car(k) > 0$ est due (avant la cohomologie \'etale) \`a Igusa, qui a utilis\'e des r\'esultats profonds sur la structure du groupe fondamental mod\'er\'e d'une courbe alg\'ebrique moins quelques points.
\end{remarque}

\begin{theoreme}[5.4]
Soient $k$ un corps alg\'ebriquement clos, et $X$ un $k$-sch\'ema propre. Alors il existe un nombre fini de courbes ferm\'ees int\`egres $Y_1, \dots, Y_p$ dans $X$, ayant la propri\'et\'e suivante : pour qu'une partie $A$ du sch\'ema $\Pic^\tau_{X/k}$ soit quasi-compacte, il faut et il suffit que les nombres d'intersection $\langle L, Y_i \rangle$ des $L \in A$ restent born\'es. De plus, on peut prendre pour $p$ le rang du groupe de N\'eron-Severi $\Pic_{X/k}(k)/\Pic^{0}_{X/k}(k)$.
\end{theoreme}

En effet, comme $\Pic^0_{X/k}$ est un sous-groupe ouvert quasi-compact de $\Pic^\tau_{X/k}$ (4.7), pour que $A$ soit quasi-compacte il faut et il suffit que l'ensemble correspondant de points de l'espace vectoriel $(\Pic^\tau_{X/k}(k)/\Pic^0_{X/k}(k)) \otimes \mathbb{Q}$ soit fini. Mais la $\tau$-\'equivalence est \'egale \`a l'\'equivalence num\'erique (4.6 (i) $\Leftrightarrow$ (iii)), donc on peut choisir parmi les courbes dans $X$ des courbes qui s'en vont sur une base de l'espace dual; d'o\`u l'assertion.


\newpage
\subsection*{6. Appendice : \'Etude des $\mathrm{(b)}$-faisceaux sur $P = \mathbb{P}^r_k$}
\addcontentsline{toc}{subsection}{6. Appendice : \'Etude des $\mathrm{(b)}$-faisceaux sur $P = \mathbb{P}^r_k$}

\subsubsection*{6.1.}
Soient $k$ un corps alg\'ebriquement clos, et $P = \mathbb{P}^r_k$. Soit $\cO_P(1)$ le faisceau inversible canonique sur $P$. Notons $S = k[x_0, \dots, x_r]$ l'anneau gradu\'e de $P$. Rappelons qu'un faisceau coh\'erent $F$ sur $P$ est dit \emph{$\mathrm{(b)}$-r\'egulier} (ou simplement $b$-r\'egulier) si $H^i(P, F(b-i)) = 0$ pour tout $i > 0$.

Le r\'esultat principal de cette section est :

\begin{theoreme}[6.2]
Soit $F$ un faisceau coh\'erent non nul sur $P$, $b$-r\'egulier. Soit $s \in H^0(P, F(b))$ une section telle que le sous-sch\'ema $Z(s)$ des z\'eros de $s$ soit de codimension $> 0$ (i.e.\ ne contienne aucune composante irr\'eductible de $\Supp(F)$). Alors le faisceau $F' = \Ker(F \xrightarrow{\,s\,} F(b))$ est $(b+1)$-r\'egulier.
\end{theoreme}

En effet, on prend $m = P_F(b)$ (polyn\^ome de Hilbert de $F$ en $b$), avec $c_i$ : $M = \sum c_i \binom{m+i}{i}$, de sorte qu'on a une suite exacte
\[
0 \to F' \to F \xrightarrow{\,s\,} F(b) \to G \to 0,
\]
o\`u $G$ est de longueur finie. Par 1.7, $F$ est $b$-r\'egulier; donc $H^i(F(n)) = 0$ pour $i > 0$ et $n \geq b-i$. Comme $\chi(G(n)) = \sum_i (-1)^i h^i(G(n)) = m$, on a
\[
h^0(G) = m, \quad h^i(G) = 0 \text{ pour } i > 0.
\]
De $s \neq 0$, il r\'esulte $G \neq 0$; par suite, $F'$ n'est pas contenu dans $F(-1)$; c'est-\`a-dire, la section $s$ n'est pas divisible par l'\'equation d'un hyperplan.

\subsubsection*{6.3.}
La r\'eciproque de 6.2 est \'egalement utile. Soit $F$ un faisceau coh\'erent sur $P$. Supposons qu'il existe une suite exacte
\[
0 \to F' \to F \xrightarrow{\,s\,} F(b) \to G \to 0
\]
avec $F'$ $(b+1)$-r\'egulier, $G$ de longueur finie. Alors $F$ est $b$-r\'egulier.

Supposons $q = 1$. Alors $\dim \Supp(N_*) = 0$; donc $h^1(N_*(b-1)) = 0$ et $h^0(N_*(b-1)) \leq h^0(N_*(b))$. De plus, $G = 0$ implique $h^0(F'(b)) = h^0(F(b)) - m$.

\begin{corollaire}[6.8. (Grothendieck)]
Soient $X$ projectif sur $S$ noeth\'erien, $\cO_X(1)$ tr\`es ample pour $X/S$, et $\mathcal{F}$ une famille de classes de faisceaux coh\'erents sur les fibres de $X/S$ limit\'ee. Alors il existe un entier $m$ tel que, pour toute classe $\xi \in \mathcal{F}$ et tout faisceau coh\'erent $F$ repr\'esentant $\xi$, $F$ soit $m$-r\'egulier sur la fibre correspondante.
\end{corollaire}

En effet, soit $X$ dans $X$. Quitte \`a remplacer $k$ par la cl\^oture alg\'ebrique d'une extension transcendante pure de $k$, on peut supposer $k$ infini. Alors, dans 6.7 (i), la premi\`ere assertion r\'esulte de 6.7 (i); la deuxi\`eme, de 6.10 et 1.8 (i), et de 6.2 (appliqu\'e \`a $F(n)$ avec $n \geq m$).

\begin{corollaire}[6.14. (L.~Hermann)]
Pour tout $N \geq 0$, il existe un polyn\^ome $R_N(x)$ tel que pour tout corps alg\'ebriquement clos $k$, pour tout $k$-sch\'ema projectif $X$, muni d'un faisceau tr\`es ample $\cO_X(1)$, et pour tout faisceau coh\'erent $F$ sur $X$ tel que $\chi(F(n)) = N$ pour tout $n$, on ait $h^0(F) \leq R_N(\deg(X))$.
\end{corollaire}

\newpage
\subsection*{7. Appendice : Th\'eor\`eme de l'indice de Hodge}
\addcontentsline{toc}{subsection}{7. Appendice : Th\'eor\`eme de l'indice de Hodge}

On retrouve et compl\`ete ici quelques r\'esultats ant\'erieurs par un raisonnement direct.

\begin{theoreme}[7.1]
Soit $X$ une surface irr\'eductible, propre sur un corps alg\'ebriquement clos. On suppose qu'il existe un faisceau inversible $H$ sur $X$ tel que $\langle c_1(H)^2 \rangle > 0$ (par exemple, $H$ ample). Alors un faisceau inversible $L$ sur $X$ tel que
\[
\langle c_1(L), c_1(H) \rangle = 0 \quad \text{et} \quad \langle c_1(L)^2 \rangle \geq 0,
\]
est num\'eriquement \'equivalent \`a $0$. (Par cons\'equent, la forme d'intersection sur $\NS(X) \otimes \mathbb{Q} = [\Pic(X)/\Pic^0(X)] \otimes \mathbb{Q}$, d'apr\`es 4.6, est non-d\'eg\'en\'er\'ee de type $(1, \rho-1)$).
\end{theoreme}

\subsubsection*{7.1.1.}
D'abord supposons : $X$ r\'eduit (donc, int\`egre); $H$ ample; $\langle c_1(L), c_1(H) \rangle = 0$; et $\langle c_1(L)^2 \rangle > 0$. Raisonnant par l'absurde, prenons $m > 0$ tel que $H_1 = L \otimes H^{\otimes m}$ soit ample, et prenons $n$ tel que, en posant $N = L^{\otimes n} \otimes H^{-1}$, on ait que $\langle c_1(N), c_1(H_1) \rangle = n\langle c_1(L)^2 \rangle - m\langle c_1(H)^2 \rangle$ est $> 0$. Alors on a aussi que $\langle c_1(N), c_1(H) \rangle = -\langle c_1(H)^2 \rangle$ est $< 0$, et que $\langle c_1(N)^2 \rangle = n^2\langle c_1(L)^2 \rangle + \langle c_1(H)^2 \rangle$ est $> 0$, ce qui contredit le lemme suivant.

\begin{lemme}[7.1.2]
Soit $X$ une surface int\`egre, projective sur un corps alg\'ebriquement clos. Pour un faisceau inversible $N$ sur $X$ tel que $\langle c_1(N)^2 \rangle > 0$, les conditions suivantes sont \'equivalentes :

\emph{(i)} Pour tout faisceau ample $H$, on a $\langle c_1(N), c_1(H) \rangle > 0$.

\emph{(i')} Il existe un faisceau ample $H$ tel qu'on ait $\langle c_1(N), c_1(H) \rangle > 0$.

\emph{(ii)} Il existe un entier $n > 0$ tel qu'on ait $H^0(X, N^{\otimes n}) \neq 0$.
\end{lemme}

En effet, (ii) implique qu'il existe un diviseur effectif $D$ tel que $N \simeq \cO_X(D)$, et $\langle c_1(N)^2 \rangle > 0$ implique que $D > 0$; d'o\`u (i).

Pour voir que (i') implique (ii), on peut \'evidemment supposer que $H = \cO_X(1)$ est tr\`es ample; alors il existe une courbe ferm\'ee int\`egre, localement principale $Y$ telle que $H = \cO_X(Y)$. D'apr\`es le th\'eor\`eme de Riemann, (i') implique que $H^p(N^{\otimes n}(-m)) = 0$ pour $p = 0$ et $n \geq n_0$, avec $n_0 = 2n'(\cO_Y) - 1$. Comme $H^p(N^{\otimes n}(m)) = 0$ pour $p > 0$ et $m \gg 0$; la suite exacte
\[
H^0(N^{\otimes n}(m)) \to H^0(N^{\otimes n}(m)|_Y) \to H^1(N^{\otimes n}(m-1)) \to H^1(N^{\otimes n}(m)) \to 0
\]
donne que $H^1(N^{\otimes n}) = 0$ pour $n \geq n_0$. Donc pour $n \gg 0$, on a
\[
h^0(N^{\otimes n}) = \chi(N^{\otimes n}) = \frac{1}{2}\langle c_1(N)^2 \rangle n^2 + \dots > 0,
\]
ce qui ach\`eve de prouver le lemme.

\subsubsection*{7.1.3.}
Maintenant supposons dans 7.1 qu'il existe un faisceau inversible $M$ sur $X$ tel que $\langle c_1(L), c_1(M) \rangle \neq 0$, et raisonnons par l'absurde. Rempla\c{c}ons $X$ par le sch\'ema r\'eduit sous-jacent d'une pr\'esentation de Chow de $X$, ce qui est loisible gr\^ace \`a la formule de projection X 4.4.4, et soit $\xi$ ample sur $X$. Prenons $m > 0$ tel que $H_1 = \xi^{\otimes m}$ soit ample, et tel que $\langle c_1(L), c_1(H_1) \rangle \neq 0$; prenons $p, q$ ($q \neq 0$) tels que, posant $L_1 = L^{\otimes p} \otimes H_1^{\otimes q}$, on ait que $\langle c_1(L_1), c_1(H) \rangle = p\langle c_1(L), c_1(H) \rangle + q\langle c_1(H_1), c_1(H) \rangle > 0$; alors on a aussi que $\langle c_1(L_1)^2 \rangle = p^2\langle c_1(L)^2 \rangle + q^2\langle c_1(H_1)^2 \rangle > 0$; ce qui revient au cas 7.1.1, avec $L = L_1$ et $H = H_1$.

\subsubsection*{7.1.4.}
Enfin consid\'erons le cas g\'en\'eral de 7.1. Toujours raisonnant par l'absurde, supposons que $L$ ne soit pas num\'eriquement \'equivalent \`a $0$; en outre, rempla\c{c}ons $X$ par la normalis\'ee $\tilde{X}$ de $X_{\mathrm{red}}$, de sorte que $\tilde{X}$ n'a qu'un nombre fini de points singuliers, r\'eduction justifi\'ee par la formule de projection X 4.4.4 et par 4.5 (ii). Il existe donc une courbe ferm\'ee int\`egre $Y$ sur $\tilde{X}$ telle que $\langle c_1(L), Y \rangle \neq 0$, et $Y$ est localement principale en dehors d'un ensemble fini $Z$. Soit $f : X' \to \tilde{X}$ l'\'eclatement de l'id\'eal $I$ de $\cO_{\tilde{X}}$. L'id\'eal $M = I \cdot \cO_{X'}$ est inversible, et il d\'efinit un sous-sch\'ema de $X'$, qui est la r\'eunion de deux sous-sch\'emas $Y'$, $Y''$ tels que $f(Y') = Y$ et $f(Y'') \subset Z$. Gr\^ace \`a la formule de projection X 4.4.4, on a que $\langle c_1(f^*L), c_1(M) \rangle = -\langle c_1(f^*L), Y' \rangle - \langle c_1(f^*L), Y'' \rangle$ est \'egal \`a $-\langle c_1(L), Y \rangle \neq 0$, et on peut remplacer $X, L, H$ par $X', f^*L, f^*H$; on est donc ramen\'e au cas pr\'ec\'edent, et la d\'emonstration de 7.1 est achev\'ee.

\begin{remarque}[7.2]
On ignore si une surface irr\'eductible, propre (sur un corps alg\'ebriquement clos), qui poss\`ede un faisceau inversible $H$ tel que $\langle c_1(H)^2 \rangle > 0$, est n\'ecessairement projective. On ignore \'egalement si pour toute surface irr\'eductible, propre, la forme d'intersection sur $\NS(X) \otimes \mathbb{Q}$ est non-d\'eg\'en\'er\'ee; il r\'esulte de 7.1 que lorsqu'elle est d\'eg\'en\'er\'ee, elle est n\'ecessairement n\'egative (non-d\'efinie).
\end{remarque}

\begin{corollaire}[7.3]
Sous les conditions de 7.1, soit $M$ un faisceau inversible sur $X$ qui n'est pas num\'eriquement \'equivalent \`a $0$, et soient $a = \langle c_1(M), c_1(H) \rangle$ et $h = \langle c_1(H)^2 \rangle$. Alors, $\langle c_1(M)^2 \rangle < a^2/h$.
\end{corollaire}

En effet, soit $L = M^{\otimes h} \otimes H^{\otimes -a}$; comme $\langle c_1(L), c_1(H) \rangle = 0$, d'apr\`es 7.1, $\langle c_1(L)^2 \rangle < 0$; d'o\`u l'assertion.

\begin{corollaire}[7.4]
Soit $X$ un sch\'ema irr\'eductible de dimension $r \geq 2$, qui est projectif sur un corps alg\'ebriquement clos, et soit $H$ un faisceau ample sur $X$. Alors :

\emph{(i)} Pour qu'un faisceau inversible $L$ sur $X$ soit num\'eriquement \'equivalent \`a $0$, (il faut et) il suffit que
\[
\langle c_1(L), c_1(H)^{r-1} \rangle = 0 \quad \text{et} \quad \langle c_1(L)^2, c_1(H)^{r-2} \rangle \geq 0.
\]

\emph{(ii)} Pour tout faisceau inversible $M$ sur $X$, qui n'est pas num\'eriquement \'equivalent \`a $0$, on a
\[
\langle c_1(M)^2, c_1(H)^{r-2} \rangle \cdot \langle c_1(H)^r \rangle < \langle c_1(M), c_1(H)^{r-1} \rangle^2.
\]
\end{corollaire}

En effet, (ii) r\'esulte tout de suite de (i), d'apr\`es le raisonnement de 7.3. Pour \'etablir (i), soit $Y$ une courbe ferm\'ee, int\`egre, dans $X$, on veut prouver que $\langle c_1(L), Y \rangle = 0$. Plongeons $X$ dans $\mathbb{P}^N$ au moyen de $H^{\otimes m}$ ($m \gg 0$). Soient $I = \bigoplus I_\nu$ l'id\'eal homog\`ene de $Y$, et $\nu_0$ un entier tel que $I_{\nu_0}$ engendre $I$ pour $\nu \geq \nu_0$. Alors les hypersurfaces $Z$ de degr\'e $\nu$, qui contiennent $Y$, n'ont pas de points communs en dehors de $Y$. Raisonnant par r\'ecurrence sur $r = \dim(X)$, on trouve \`a l'aide du th\'eor\`eme de Bertini $r-2$ hypersurfaces $Z_1, \dots, Z_{r-2}$ de degr\'e $\nu$, qui contiennent $Y$, telles que $X' = Z_1 \cap \dots \cap Z_{r-2} \cap X$ soit irr\'eductible. Comme $\langle c_1(L), Y \rangle = \langle c_1(L|_{X'}), Y \rangle$, il suffit de prouver que $L|_{X'}$ est num\'eriquement \'equivalent \`a $0$ sur $X'$. Or par la formule de projection, on a que $\langle c_1(L|_{X'}), c_1(H|_{X'}) \rangle = 0$ et $\langle c_1(L|_{X'})^2 \rangle \geq 0$; donc l'assertion r\'esulte de 7.1.

\begin{corollaire}[7.5]
Soient $k$ un corps alg\'ebriquement clos, $Y$ un $k$-sch\'ema projectif, muni d'un faisceau ample $H$, et $X$ le sch\'ema des z\'eros d'une section de $H$. On suppose que pour toute composante irr\'eductible $Y^*$ de $Y$, les composantes irr\'eductibles de $X \cap Y^* = V(\omega|_{Y^*})$ sont de dimension $> 2$. Alors le morphisme $\NS(Y) \to \NS(X)$ est injectif.
\end{corollaire}

En effet, comme on peut \'evidemment supposer $Y$ int\`egre, l'assertion r\'esulte de 7.4 (i), gr\^ace \`a la formule de projection.

\subsubsection*{Bibliographie}
\begin{enumerate}[label={[\arabic*]}]
\item[{[1]}] Abhyankar S. : R\'esolution of singularities of arithmetical surfaces, in : Arithmetical Algebraic Geometry, Purdue University, Harper \& Row, New York 1965.
\item[{[2]}] Mumford D. : Pathologies III, Amer.\ J.\ Math.\ 1967 p.\ 94--103.
\item[{[3]}] Grothendieck A. et Dieudonn\'e J. : EGA, \'El\'ements de G\'eom\'etrie Alg\'ebrique, Pub.\ Math.\ Inst.\ des Hautes \'Etudes Sci., 1960.
\item[{[4]}] Grothendieck A. : FGA, Fondements de la g\'eom\'etrie alg\'ebrique, Extraits du S\'eminaire Bourbaki 1957--1962.
\item[{[5]}] Grothendieck A. et alii, SGA, S\'eminaire de g\'eom\'etrie alg\'ebrique, 1960--66, I.H.\'E.S.
\item[{[6]}] Jussila O. : X, Formalisme des intersections sur les sch\'emas alg\'ebriques propres, Expos\'e X dans ce s\'eminaire.
\item[{[7]}] Raynaud M. : XII, Un th\'eor\`eme de repr\'esentabilit\'e relative sur le foncteur de Picard, Expos\'e XII dans ce s\'eminaire.
\end{enumerate}


% ============================================================
\newpage
\section*{Expos\'e XIV}
\addcontentsline{toc}{section}{Expos\'e XIV --- Probl\`emes ouverts en th\'eorie des intersections}

\begin{center}
\Large\textbf{PROBL\`EMES OUVERTS EN TH\'EORIE DES INTERSECTIONS}

\medskip
\normalsize\textit{par A.\ Grothendieck}
\end{center}

\bigskip

Dans le pr\'esent expos\'e nous passons en revue quelques uns des probl\`emes li\'es aux questions rencontr\'ees dans ce s\'eminaire, et qui m\'eriteraient d'\^etre \'etudi\'es. Nous laissons de c\^ot\'e les probl\`emes li\'es \`a la th\'eorie locale des intersections \`a la Serre [14], que nous n'avons pas eu \`a utiliser dans le S\'eminaire, et renvoyons notamment \`a loc.\ cit.\ pour les deux conjectures fondamentales qui restent ouvertes dans cette direction (loc.\ cit.\ Chap.\ V \S\ 4). Nous passons sous silence aussi les probl\`emes, fort importants pour les applications arithm\'etiques, li\'es \`a la structure des groupes de classes de cycles, sur un sch\'ema alg\'ebrique projectif et lisse notamment; nous renvoyons pour la discussion de ces probl\`emes \`a S.\ KLEIMAN, \textit{Algebraic cycles and the Weil conjectures} dans Dix expos\'es sur la cohomologie des sch\'emas, (Masson--North Holland Publishing Cie). Dans le n\textsuperscript{o} 4 du pr\'esent expos\'e, nous donnons \'egalement quelques compl\'ements au S\'eminaire, concernant notamment l'anneau de Chow, qui n'avaient trouv\'e leur place dans les expos\'es pr\'ec\'edents.

\subsection*{1. Op\'erations $\lambda$ dans la cat\'egorie d\'eriv\'ee $D(X)$}
\addcontentsline{toc}{subsection}{1. Op\'erations $\lambda$ dans la cat\'egorie d\'eriv\'ee $D(X)$}

Soit $X$ un topos localement annel\'e, et soit $K^\bullet(X)$ l'anneau des classes de complexes parfaits sur $X$ (Exp.\ IV). On aimerait d\'efinir dans $K^\bullet(X)$ une structure de $\lambda$-anneau (Exp.\ V), fonctorielle en $X$, telle que l'homomorphisme naturel
\[
K^\bullet(X) \to K^\bullet(X)
\]
(o\`u le premier membre est l'anneau des classes de Modules localement libres sur $X$, muni de la $\lambda$-structure d\'efinie dans Exp.\ V) soit un $\lambda$-homomorphisme. Dans ce but, il y aurait lieu de d\'efinir des foncteurs (non additifs en g\'en\'eral) de ``puissance ext\'erieure''
\[
\Lambda^i : D(X) \to D(X),
\]
transformant complexes pseudo-coh\'erents (resp.\ parfaits) en complexes pseudo-coh\'erents (resp.\ parfaits), et ayant les propri\'et\'es formelles n\'ecessaires, vis \`a vis des triangles exacts, pour d\'efinir une $\lambda$-structure sur $K^\bullet(X)$ par la formule
\[
\lambda^i(\cl(L_\bullet)) = \cl(\Lambda^i(L_\bullet)). \tag{1.1}
\]

Lorsque $X$ est de caract\'eristique nulle, i.e.\ son faisceau structural est un faisceau de $\mathbb{Q}$-alg\`ebres, le probl\`eme se r\'esoud simplement en prenant
\[
\Lambda^i(L_\bullet) = (L_\bullet^{\otimes i})^{\mathrm{alt}}, \tag{1.2}
\]
o\`u l'on consid\`ere la puissance tensorielle $i$-\`eme (au sens des cat\'egories d\'eriv\'ees) de $L_\bullet$ comme \'etant un objet de $D(\cO_X[\mathfrak{S}_i])$, $\cO_X[\mathfrak{S}_i]$ \'etant l'alg\`ebre du groupe sym\'etrique $\mathfrak{S}_i$ \`a co\'efficients dans $\cO_X$, et o\`u l'exposant ``alt'' d\'esigne le passage \`a la ``partie altern\'ee sous l'action de $\mathfrak{S}_i$'' (op\'eration qui a un sens gr\^ace \`a l'hypoth\`ese de caract\'eristique nulle). Une formule de ce type \'etait utilis\'ee d\'ej\`a dans [8] (cit\'e par la suite [RRR]), pour la premi\`ere d\'emonstration du th\'eor\`eme de Riemann-Roch.

Sans hypoth\`ese de caract\'eristique, lorsque $L_\bullet$ est nul en degr\'es $> 0$, alors par Dold--Puppe [6] $L_\bullet$ est canoniquement homotope au complexe associ\'e \`a un Module semi-simplicial $L_\bullet'$ canoniquement associ\'e \`a $L_\bullet$, et un candidat assez naturel pour $\Lambda^i(L_\bullet)$ (pour $L_\bullet$ plat, disons) serait le complexe associ\'e au Module semi-simplicial $\Lambda^i L_\bullet'$ (dont la composante de degr\'e $n$ est $\Lambda^i L_n'$). Mais pour \'etendre la d\'efinition \`a des complexes de $D(X)$ qui ne sont plus suppos\'es nuls en degr\'es $> 0$, il faudrait \'etudier le comportement de l'op\'eration pr\'ec\'edente par rapport \`a l'op\'eration de translation $L_\bullet \mapsto L_\bullet[1]$ (ou suspension, dans la terminologie de [6]); malheureusement, il ne semble pas qu'on ait un isomorphisme $\Lambda^i(L_\bullet[1]) \simeq \Lambda^i(L_\bullet)[i]$ qui permettrait de faire de fa\c{c}on triviale l'extension d\'esir\'ee.

\subsection*{2. La formule de Riemann-Roch sans hypoth\`eses projectives}
\addcontentsline{toc}{subsection}{2. La formule de Riemann-Roch sans hypoth\`eses projectives}

\subsubsection*{2.1.}
Une fois en possession d'une $\lambda$-structure naturelle sur les $K^\bullet(X)$, il est possible de formuler le th\'eor\`eme de Riemann-Roch, sous la forme envisag\'ee dans Exp.\ VIII, pour n'importe quel morphisme $f : X \to Y$ propre et d'intersection compl\'ete, (et m\^eme, plus g\'en\'eralement, dans le cas d'un sch\'ema relatif $X$ sur un topos localement annel\'e $Y$, $X$ \'etant encore soumis aux conditions relatives correspondantes). La classe de Todd de $f$ se d\'efinit gr\^ace \`a la classe dans $K^\bullet(X)$ du complexe cotangent relatif $\mathbb{L}_\bullet$ pour $f$, construit dans Exp.\ 0 4.4, lequel complexe est ici parfait gr\^ace \`a l'hypoth\`ese locale faite sur $f$. Il semble clair que la d\'emonstration de la formule de Riemann-Roch dans ce cas g\'en\'eral demandera l'introduction d'id\'ees essentiellement nouvelles, et il est possible que ces m\^emes id\'ees pourraient donner la clef d'une d\'emonstration de la formule analogue dans le cas analytique complexe (ou rigide-analytique !). Rappelons que dans le cas analytique complexe, seul le cas d'un faisceau localement libre $F$ sur une vari\'et\'e non singuli\`ere compacte $X$, au-dessus d'un $Y$ ponctuel, est actuellement connu (gr\^ace \`a la formule de l'index de Atiyah--Singer), le cas o\`u $F$ est seulement suppos\'e coh\'erent semblant \'echapper encore aux techniques connues jusqu'\`a pr\'esent.

\subsubsection*{2.2.}
Comme ingr\'edient essentiel de la formule de Riemann-Roch, rappelons aussi pour m\'emoire le probl\`eme de finitude soulev\'e dans Exp.\ III : si $f : X \to Y$ est un morphisme propre, et $L_\bullet$ un complexe sur $X$ qui est pseudo-coh\'erent relativement \`a $Y$, alors $\RR f_*(L_\bullet)$ est-il pseudo-coh\'erent ? Une d\'emonstration g\'en\'erale de cette conjecture mettra probablement en \oe{}uvre des id\'ees alg\'ebriques qui permettront de donner une d\'emonstration du ``th\'eor\`eme'' de finitude de Grauert [7] sous des conditions notablement plus g\'en\'erales que dans loc.\ cit., en se pla\c{c}ant d'embl\'ee sur des espaces (ou topos) de base topologiquement annel\'es en $\mathbb{C}$-alg\`ebres $Y$ plus g\'en\'eraux que des esp\`eces analytiques (englobant le cas o\`u $Y$ serait un espace localement compact avec le faisceau des fonctions continues complexes, ou une vari\'et\'e $C^\infty$ avec le faisceau des fonctions complexes $C^\infty$, ou un espace analytique r\'eel, etc).

\subsubsection*{2.3.}
Mentionnons enfin qu'il y aurait lieu de remplacer dans ces questions l'hypoth\`ese de propret\'e sur le morphisme $f$ par l'hypoth\`ese que les faisceaux de cohomologie $H^i(L_\bullet)$ des complexes envisag\'es sur $X$ aient leurs groupes propres sur $Y$.


\subsection*{3. Formule de Riemann-Roch ``sans denominateurs'' pour une immersion}
\addcontentsline{toc}{subsection}{3. Formule de Riemann-Roch ``sans denominateurs'' pour une immersion}

\subsubsection*{3.1.}
Soit $f : Y \to X$ une immersion fermee reguliere (Exp.\ VII), de faisceau conormal $N$. Supposons pour simplifier que $X$ admet un faisceau inversible ample. La formule de Riemann-Roch pour $y \in K^{\bullet}(Y)$ :
\[
\ch(f_{*}(y)) = f_{*}(\ch(y) \cdot \Td(-N))
\]
jointe a la formule generale (Exp.\ VII)
\[
f_{*}(y) \cdot \xi = f_{*}(y \cdot f^{*}(\xi)) \quad (\xi \in K^{\bullet}(X)),
\]
permet d'exprimer les classes de Chern de $f_{*}(y)$ dans $\operatorname{Gr}^{i}(X)$ sous la forme
\[
c_i(f_{*}(y)) = P_i(c_j(y), c_k(N)), \tag{3.1}
\]
ou $P_i$ est un polynome universel en les $c_j(y)$, $c_k(N)$, polynome qui a priori est a coefficients rationnels. En fait, on constate [RRR] que ces polynomes sont a coefficients entiers. Cela souleve la question de la validite d'une formule de Riemann-Roch ``sans denominateurs''. En fait, une telle formule pour $c_i(f_{*}(y))$ ne peut etre cherchee en general dans $\operatorname{Gr}^{\bullet}(X)$ lui-meme, vu que l'homomorphisme $f_{!} : K^{\bullet}(Y) \to K^{\bullet}(X)$ n'est pas compatible lui-meme avec les $\gamma$-filtrations modulo decalage par $d$ (4.6), mais l'est seulement a torsion pres; de sorte qu'il n'est pas possible de definir un homomorphisme $f_{*} : \operatorname{Gr}^{\bullet}(Y) \to \operatorname{Gr}^{\bullet+d}(X)$ donnant l'homomorphisme naturel $f_{!} : \operatorname{Gr}^{\bullet}(Y)_{\mathbb{Q}} \to \operatorname{Gr}^{\bullet+d}(X)_{\mathbb{Q}}$, et permettant de donner un sens a la formule (3.1) comme une formule dans $\operatorname{Gr}^{\bullet}(X)$, et non seulement dans $\operatorname{Gr}^{\bullet}(X)_{\mathbb{Q}}$. Mais lorsque $X$ est suppose regulier, donc $K^{\bullet}(X) \simeq K_{\bullet}(X)$ (IV.2.5), on peut munir $K^{\bullet}(X)$ de la filtration topologique par la codimension des supports, heritee de la filtration correspondante de $K_{\bullet}(X)$. Il est plausible que cette filtration est compatible avec la structure d'anneau, fait qu'on peut prouver tout au moins [RRR] lorsque $X$ est lisse sur un corps $k$, en utilisant la theorie de l'anneau de Chow [4]. Quand il en est ainsi, le gradue associe
\[
\operatorname{Gr}_{\mathrm{top}}^{\bullet}(X)
\]
est un anneau gradue, et supposant egalement $Y$ regulier, on a fait ce qu'il fallait pour avoir un homomorphisme naturel de degre $d$ de groupes gradues
\[
f_{*} : \operatorname{Gr}_{\mathrm{top}}^{\bullet}(Y) \longrightarrow \operatorname{Gr}_{\mathrm{top}}^{\bullet+d}(X).
\]
Comme la $\gamma$-filtration de $K^{\bullet}(X)$ est plus fine que la filtration topologique par la codimension qu'on vient de considerer (X 1.3.3), on trouve un homomorphisme canonique de groupes gradues (et meme d'anneaux gradues)
\[
\operatorname{Gr}^{\bullet}(X) \longrightarrow \operatorname{Gr}_{\mathrm{top}}^{\bullet}(X)
\]
qui est d'ailleurs un isomorphisme mod torsion (Exp.\ VIII). Ceci permet, comme dans [RRR], de considerer une theorie des classes de Chern sur $X$ a valeurs dans $\operatorname{Gr}_{\mathrm{top}}^{\bullet}(X)$, et de meme sur $Y$, et de donner ainsi un sens a la formule (3.1) comme une formule dans $\operatorname{Gr}_{\mathrm{top}}^{\bullet}(X)$. Sous cette forme, cette formule est prouvee dans [RRR] lorsque $X$ est de caracteristique nulle, tout au moins si $X$ est lisse sur un corps (hypothese servant uniquement a assurer que la filtration topologique de $K^{\bullet}(X)$ est compatible avec la structure d'anneau). L'hypothese de caracteristique nulle est utilisee essentiellement pour disposer de formules du type (1.1), (1.2), et la meme demonstration essentiellement devrait s'appliquer sans hypothese de caracteristique, une fois resolue dans un sens affirmatif la question soulevee au n\textsuperscript{o} 1.

\subsubsection*{3.2.}
Lorsque $X$ est lisse sur un corps, on dispose aussi de l'anneau de Chow $A(X)$, et d'une theorie des classes de Chern sur $X$ a valeurs dans $A(X)$ [9], plus fine que la theorie a valeurs dans $\operatorname{Gr}_{\mathrm{top}}^{\bullet}(X)$ qu'on vient d'introduire, ces dernieres classes de Chern etant images des classes de Chern style Chow par l'homomorphisme canonique surjectif (cf.\ 4.1)
\[
\phi : A(X) \longrightarrow \operatorname{Gr}_{\mathrm{top}}^{\bullet}(X).
\]
Cet homomorphisme est un isomorphisme mod.\ torsion, mais pas un isomorphisme en general (4.7). Ceci dit, si $X$ et $Y$ sont lisses, l'inclusion $f : Y \hookrightarrow X$ definit encore un homomorphisme de degre $d$ de groupes gradues
\[
f_{*} : A^{\bullet}(Y) \longrightarrow A^{\bullet+d}(X),
\]
ce qui permet de donner un sens a la formule (3.1) comme une formule dans $A(X)$; cette formule se reduit d'ailleurs a $0 = 0$ pour $i < d$, et ne presente donc d'interet que pour $i > d$. Meme pour $i = d = 2$ cette formule n'est pas prouvee a l'heure actuelle (cf.\ 4.3), ni la formule moins precise qui s'en deduit en remplacant l'equivalence rationelle des cycles algebriques par l'equivalence algebrique. Contrairement a la formule analogue dans $\operatorname{Gr}_{\mathrm{top}}^{\bullet}(X)$, il ne semble pas y avoir de raison bien convaincante pour supposer que la formule sans denominateurs doive etre vraie dans $A(X)$. Une bonne facon de la tester, vue notre ignorance quasi totale sur la nature des groupes $A^{i}(X)$ pour $i > 2$, semble etre de tester les formules qu'on peut en deduire dans des $A^{2}(Z) = \Pic(Z)$, par exemple en transformant les deux membres de (3.1) par $g_{*} : A(X) \to A(Z)$, ou $g : X \to Z$ est un morphisme propre dans un $Z$ quasi-projectif lisse de dimension $\dim X - i + 1$.

\subsubsection*{3.3.}
Enfin, pour donner un sens a (3.1), on peut egalement utiliser la theorie cohomologique des classes de Chern (SGA 5 VII), chaque fois qu'on dispose d'homomorphismes de Gysin
\[
f_{*} : H^{\bullet-2d}(Y, \mathbb{Z}_{\ell}(\bullet-d)) \longrightarrow H^{\bullet}(X, \mathbb{Z}_{\ell}(\bullet))
\]
(cf.\ n\textsuperscript{o} 6), ce qui est le cas tout au moins si $Y$ et $X$ sont tous deux lisses sur un schema $S$ (grace au theoreme de purete cohomologique relatif SGA 4 XVI 3.7), ce, conjecturalement, (et apres le theoreme de purete cohomologique absolu, demontre tout au moins pour $X$ excellent de caracteristique nulle (SGA 4 XIX 3.2)). Lorsque $X$ et $Y$ sont lisses sur le corps des complexes, on peut egalement utiliser la classe de Chern entiere sur l'espace localement compact $X(\CC)$ des points complexes de $X$. La formule (3.1) est etablie dans ce cas par Atiyah et Hirzebruch [1], qui se placent plus generalement dans le cas d'une immersion de varietes analytiques complexes (en travaillant avec le $K^{\bullet}$ ``naif''). Il en resulte que la formule (3.1) $l$-adique est valable chaque fois que $Y$ et $X$ sont lisses sur un corps algebriquement clos et de caracteristique nulle, en utilisant le ``principe de Lefschetz'' et les resultats generaux de SGA 4. Cela rend donc la validite de la formule (3.1) cohomologique extremement plausible; mais on notera qu'elle n'est pas demontree a l'heure actuelle, meme si $X$ et $Y$ sont lisses sur un corps $k$ de caracteristique nulle $k$, lorsque $k$ n'est pas algebriquement clos.

\subsubsection*{3.4.}
On peut egalement travailler avec la theorie cohomologique des classes de Chern, style De Rham ou style Hodge [13]. Le probleme de la validite de (3.1) dans ce cas ne se pose que si $X$ n'est pas de caracteristique nulle, puisque l'on sait de toutes facons (n\textsuperscript{o} 6) que la formule en question est vraie modulo torsion. Pour les schemas lisses quasi-projectifs sur un corps $k$, la formule (3.1) version De Rham ou Hodge serait d'ailleurs consequence de la formule (3.1) version classes de cycles mod.\ equivalence algebrique, de sorte que la version cohomologique De Rham ou Hodge donne une autre facon abordable de tester la validite de (3.1) version ``classes de cycles''.

\subsubsection*{3.5.}
La demonstration [1] de la version cohomologique transcendante de (3.1) est de nature essentiellement transcendante, en faisant usage de voisinages tubulaires de $Y$ dans $X$. Cela suggere que la demonstration d'une formule du type (3.1) pourrait etre liee a la consideration systematique de groupes $K^{\bullet}$ et de groupes de cohomologie a support dans $Y$, consideration qui pourrait etre un moyen technique suffisant de ``localisation autour de $Y$''.


\subsection*{4. Relations entre $K^{\bullet}(X)$ et l'anneau de Chow $A(X)$}
\addcontentsline{toc}{subsection}{4. Relations entre $K^{\bullet}(X)$ et l'anneau de Chow $A(X)$}

\subsubsection*{4.1.}
Supposons pour simplifier que $X$ soit un schema lisse et quasi-projectif sur un corps $k$ (cf.\ n\textsuperscript{o} 8 pour un cas plus general), de sorte qu'on dispose de l'anneau de Chow $A(X)$ des classes de cycles a equivalence rationnelle pres [4]. Comme variantes de cet anneau, definis d'ailleurs sous des considerations sensiblement plus generales, on a les anneaux $\operatorname{Gr}^{\bullet}(X)$ (gradue associe a la $\gamma$-filtration de $K^{\bullet}(X)$) et $\operatorname{Gr}_{\mathrm{top}}^{\bullet}(X)$ (gradue associe a la filtration de $K^{\bullet}(X) = K_{\bullet}(X)$ par la codimension des supports). On a des theories des classes de Chern a valeurs dans chacun de ces trois anneaux (cf.\ [9] pour le cas de $A(X)$, et ce seminaire pour $\operatorname{Gr}^{\bullet}(X)$, qui s'envoie dans $\operatorname{Gr}_{\mathrm{top}}^{\bullet}(X)$). On a egalement une theorie cohomologique ($l$-adique, disons, ou $l$ est premier a la caracteristique de $k$) des classes de Chern (SGA 5 VII) :
\[
c_i(E) \in H^{2i}(X, \mathbb{Z}_l(i))
\]
(ou $E$ est un fibre vectoriel sur $X$, ou encore un element quelconque de $K^{\bullet}(X)$). Ces theories cohomologiques sont reliees a l'aide des homomorphismes canoniques (compatibles avec les structures multiplicatives)
\[
\xymatrix@R=0.5pc{
& A^{\bullet}(X) \ar[dd]^{\phi} \ar[rdd]^{\psi} & \\
\operatorname{Gr}^{\bullet}(X) \ar[rd]^{j} \ar[rrd]_(0.3){\chi} & & \\
& \operatorname{Gr}_{\mathrm{top}}^{\bullet}(X) & H^{2\bullet}(X, \mathbb{Z}_l(\bullet))
} \tag{4.1}
\]
ou on a pose
\[
H^{2\bullet}(X, \mathbb{Z}_l(\bullet)) = \bigoplus_i H^{2i}(X, \mathbb{Z}_l(i)).
\]
L'homomorphisme $j$ provient du fait que la $\gamma$-filtration de $K^{\bullet}(X)$ est plus fine que la filtration topologique par la codimension du support (Exp.\ X), et c'est un isomorphisme modulo torsion (Exp.\ VII 4.11). L'homomorphisme $\phi$ est defini par exemple dans [RRR], et associe a la classe du cycle de codimension $i$ associe a un sous-schema ferme $Y$ de $X$ de codimension $i$, l'element $\cl(\cO_Y)$ de $\Fil_{\mathrm{top}}^i(K_{\bullet}(X)) \mod \Fil_{\mathrm{top}}^{i+1}(K_{\bullet}(X))$. C'est un homomorphisme surjectif, dont nous verrons au n\textsuperscript{o} 4.2 que c'est un isomorphisme modulo torsion. Ainsi, modulo torsion les trois anneaux $A(X)$, $\operatorname{Gr}_{\mathrm{top}}^{\bullet}(X)$ et $\operatorname{Gr}^{\bullet}(X)$ sont canoniquement isomorphes, donc a posteriori les formules de Riemann-Roch exprimees a l'aide de l'un ou l'autre de ces anneaux (tensorises par $\mathbb{Q}$) sont equivalentes. L'anneau gradue $\operatorname{Gr}^{\bullet}(X)$, surtout utilise dans ce seminaire comme substitut a l'anneau de Chow, a l'avantage appreciable d'etre defini pour tout topos localement annele, et en particulier pour tout schema (sans hypothese de regularite ni d'existence de Module inversible ample), et d'acquerir un caractere covariant par rapport aux morphismes propres (lisses ou non). Il a par contre le desavantage (sur $A(X)$, $\operatorname{Gr}_{\mathrm{top}}^{\bullet}(X)$) de n'acquerir ce caractere qu'une fois tensorise par $\mathbb{Q}$. Autre avantage de $A(X)$ sur ses concurrents $\operatorname{Gr}_{\mathrm{top}}^{\bullet}(X)$ et $\operatorname{Gr}^{\bullet}(X)$, c'est qu'on dispose d'un homomorphisme de $A(X)$ dans $H^{2\bullet}(X, \mathbb{Z}_l(\bullet))$, permettant d'interpreter les classes de Chern cohomologiques comme images des classes de Chern a valeurs dans $A(X)$. Cet homomorphisme en general ne se factorise pas par $\operatorname{Gr}_{\mathrm{top}}^{\bullet}(X)$ (4.7), et il est pour le moins douteux qu'il existe en general un homomorphisme de $\operatorname{Gr}^{\bullet}(X)$ dans $H^{2\bullet}(X, \mathbb{Z}_l(\bullet))$ qui transforme classes de Chern en classes de Chern (cf.\ n\textsuperscript{o} 5). Bien entendu, ces difficultes disparaissent lorsqu'on tensorise par $\mathbb{Q}$, et on trouve un homomorphisme
\[
\operatorname{Gr}^{\bullet}(X)_{\mathbb{Q}} \simeq \operatorname{Gr}_{\mathrm{top}}^{\bullet}(X)_{\mathbb{Q}} \simeq A^{\bullet}(X)_{\mathbb{Q}} \longrightarrow H^{2\bullet}(X, \mathbb{Q}_l(\bullet)) \tag{4.2}
\]
qui transforme classes de Chern en classes de Chern.

\subsubsection*{4.2.}
Pour etablir que l'homomorphisme surjectif $\phi$ de (4.1) est un isomorphisme modulo torsion, on utilise l'homomorphisme de Chern (deduit de la theorie des classes de Chern [9])
\[
K_{\bullet}(X) \longrightarrow \Ch(A(X))
\]
(ou le deuxieme membre est l'anneau de Chern de $A(X)$ defini dans Exp.\ V), qui est compatible comme on voit aisement avec la filtration topologique de $K^{\bullet}(X)$ et la filtration naturelle de $\Ch(A(X))$, et donne par passage aux gradues associes un homomorphisme de groupes gradues
\[
\psi : \operatorname{Gr}_{\mathrm{top}}^{\bullet}(X) \longrightarrow A(X). \tag{4.3}
\]
Ceci pose, on va prouver les formules
\[
\psi \circ \phi = (-1)^{i-1}(i-1)! \cdot x \mod.\ \text{torsion} \quad \text{pour } x \in A^{i}(X) \tag{4.4}
\]
\[
\phi \circ \psi = (-1)^{i-1}(i-1)! \cdot x \mod.\ \text{torsion} \quad \text{pour } x \in \operatorname{Gr}_{\mathrm{top}}^{i}(X),
\]
dont la premiere implique d'ailleurs la seconde, grace au fait que $\phi$ est surjectif. Ces formules impliquent evidemment que $\phi$ et $\psi$ sont des isomorphismes modulo torsion, et qu'a un facteur $(-1)^{i-1}(i-1)!$ pres en degre $i$, ils sont inverses l'un de l'autre.

Pour prouver la premiere formule (4.4), on peut evidemment supposer que $x$ est defini par un sous-schema ferme irreductible $Y$ de $X$ de codimension $i$, auquel cas la formule equivaut a
\[
\cl(\cO_Y) = (-1)^{i-1}(i-1)! \, \cl(Y) \mod \text{ torsion} \tag{4.5}
\]
ou dans le premier membre $\cO_Y$ est considere comme un faisceau coherent sur $X$, definissant donc un element de $K^{\bullet}(X)$, et dans le deuxieme membre $\cl(Y)$ designe la classe de $Y$ dans $A^{i}(X)$. Comme $A^{i}(X)$ ne change pas en enlevant de $X$ une partie fermee de codimension $> i$, on est ramene au cas ou $Y$ est regulier, donc lisse sur le corps de base $k$ si ce dernier est parfait. Mais dans ce dernier cas, la formule (4.5) n'est autre que la formule de Riemann-Roch pour l'inclusion $Y \hookrightarrow X$ et l'element $1$ de $K^{\bullet}(Y)$, formule ecrite en utilisant la theorie des classes de Chern a valeurs dans l'anneau de Chow $A(X)$. Cette formule est etablie dans [2] sous les conditions envisagees ici, par une methode tres voisine de celle utilisee dans ce seminaire pour la formule analogue utilisant $\operatorname{Gr}^{\bullet}(X)$. Cela etablit donc, du moins pour $k$ parfait, la formule (4.4) et par suite le fait que $\phi$ soit un isomorphisme mod.\ torsion (qui entraine a son tour, a posteriori, que la formule de Riemann-Roch ``avec denominateurs'', pour un morphisme propre de schemas algebriques lisses sur $k$, est essentiellement la meme, qu'on l'ecrive en utilisant $A(X)$ comme dans [2], ou en utilisant $\operatorname{Gr}^{\bullet}(X)$ comme dans ce seminaire !). Le cas ou on ne fait pas d'hypothese sur $k$ se ramene aisement au cas ou $k$ est algebriquement clos, en utilisant le fait que pour une extension finie $k'$ de degre $n$ de $k$, le noyau de l'application naturelle $A(X) \to A(X \otimes_k k')$ est annule par $n$ (comme on voit en utilisant l'homomorphisme trace $f_{*} : A(X \otimes_k k') \to A(X)$, satisfaisant $f_{*}f^{*} = n \cdot \id$).

\subsubsection*{4.3.}
Il semble assez plausible que les formules (4.4) soient vraies telles quelles et non seulement modulo torsion, ou ce qui revient au meme, que la formule (4.5) soit vraie non seulement modulo torsion. Cette question est un cas tres particulier de celle de la validite dans $A(X)$ de la ``formule de Riemann-Roch sans denominateurs'' (3.1), pour une immersion $Y \to X$, (savoir le cas ou $i = d$). On peut considerer la validite de ce cas particulier comme assez plausible. En effet, nous avons signale deja au n\textsuperscript{o} 3 que lorsque $k$ est de caracteristique nulle, la formule de Riemann-Roch sans denominateurs est valable tout au moins dans $\operatorname{Gr}_{\mathrm{top}}^{\bullet}(X)$, ce qui implique dans ce cas que les deux membres de (4.5) ont meme image dans $\operatorname{Gr}_{\mathrm{top}}^{i}(X)$, ou ce qui revient au meme, que la deuxieme des formules (4.4) est valable; ce resultat s'etendra de lui-meme au cas de la caracteristique quelconque, une fois resolue de facon satisfaisante la question soulevee au n\textsuperscript{o} 1.

Comme support pour la formule (4.5) dans $A(X)$ lui-meme, signalons le fait qu'il est clair a priori que la restriction du premier membre a $X - Y$ est nulle, donc, $Y$ etant suppose irreductible, le premier membre doit etre a priori un multiple entier de $\cl(Y)$ :
\[
\cl(\cO_Y) = a(Y) \cdot \cl(Y).
\]
Cela prouve deja la formule exacta (4.5) lorsque $\cl(Y)$ n'est pas un element de torsion de $A^{i}(X)$. Il est immediat que tel est le cas si $X$ est projectif sur $k$ et si $\cl(Y) \cdot \xi^{r-i} > 0$ pour une polarisation $\xi$ de $X$, alors $\deg(Y \cdot \xi^{r-i}) > 0$; cela prouve la formule voulue dans ce cas. Il en resulte que cette formule est encore valable chaque fois que $X$ est isomorphe a un sous-schema ouvert d'un schema projectif lisse sur $k$, ce qui sera le cas lorsqu'on dispose d'une solution convenable du probleme de resolution des singularites. Il en est ainsi en particulier si $k$ est de caracteristique nulle, grace a Hironaka [11].

La validite de la formule exacte (4.5) donc de (4.4) impliquerait que le noyau de $\phi : A^{i}(X) \to \operatorname{Gr}_{\mathrm{top}}^{i}(X)$ est annule par $(i-1)!$, et en particulier que $\phi : A^{2}(X) \to \operatorname{Gr}_{\mathrm{top}}^{2}(X)$ est un isomorphisme. (Que $\phi^1$ soit un isomorphisme resulte deja de X 5.3.2). Sauf erreur, le fait que $\phi$ soit un isomorphisme n'est pas prouve a l'heure actuelle, meme dans le cas particulier ou $\dim X = 3$. C'est hativement en tous cas que le redacteur avait affirme la validite de (4.4) et (4.5) dans [9, n\textsuperscript{o} 4, 3\textsuperscript{o}].

\subsubsection*{4.4.}
A l'heure actuelle, le principal avantage de $\operatorname{Gr}_{\mathrm{top}}^{\bullet}(X)$ (ou de $\operatorname{Gr}^{\bullet}(X)$) sur l'anneau de Chow $A(X)$ est le fait qu'il se prete mieux aux calculs explicites de certaines formules d'intersection entre cycles, dans le cas d'intersections ``excedentaires'' i.e.\ de dimension trop grande. La definition de la classe d'intersections, a l'aide du lemme de Chow permettant de bouger suffisamment un cycle dans sa classe, se prete en effet malaisement au calcul. Un exemple frappant a cet egard est la formule fondamentale suivante [Exp.\ VII 2.7], pour une immersion fermee $f : Y \to X$ de schemas quasi-projectifs lisses sur $k$ :
\[
f_{*}(y) = y \cdot c_d(N), \tag{4.6}
\]
formule valable dans $\operatorname{Gr}_{\mathrm{top}}^{\bullet}(X)$, $y$ etant dans $\operatorname{Gr}^{\bullet}(Y)$, $N$ le faisceau conormal de $Y$ dans $X$, suppose partout de rang $d$. Dans loc.\ cit., cette formule resulte d'un calcul tres simple de faisceaux $\Tor$. Grace a ce qu'on a vu dans 4.2, la formule precedente implique la meme formule dans $A(Y)$, mais seulement modulo torsion. Il est assez extraordinaire que la formule (4.6) dans $A(Y)$ lui-meme ne soit pas prouvee a l'heure actuelle\footnote{Ajoute en avril 1969 : D.\ Mumford nous signale une demonstration de (4.6) dans l'anneau de Chow (sans negliger la torsion), qu'il a obtenue dans un seminaire a Harvard en 1959. Le lecteur trouvera la demonstration de ce beau resultat dans SGA 5 VII 9.}, ni meme le cas particulier obtenu en faisant $y = 1$, donnant la self-intersection de $Y$ :
\[
f_{*}(1) = c_d(N), \tag{4.7}
\]
meme si $k$ est le corps des complexes !

On se heurte a un probleme analogue lorsqu'on tente de calculer l'anneau de Chow de la variete eclatee $X'$ de $X$ le long de $Y$, en calquant la methode suivie pour le calcul de $\operatorname{Gr}^{\bullet}(X')$, (Exp.\ VII 4.8), et qui marche egalement (sans tensoriser par $\mathbb{Q}$) pour le calcul de $\operatorname{Gr}_{\mathrm{top}}^{\bullet}(X')$. La formule-cle pour cette determination (avec les notations de loc.\ cit.) est
\[
c_i(\xi) = j_{*}(\pi^{*}(c_{i-d}(F))) \tag{4.8}
\]
formule etablie a l'heure actuelle dans $\operatorname{Gr}_{\mathrm{top}}^{\bullet}(X')$ (pour $y \in \operatorname{Gr}_{\mathrm{top}}^{\bullet}(Y)$) grace a un calcul de $\Tor$ (Exp.\ VII 3.4), et dont la validite dans $A(X')$ (pour $y \in A(Y)$) resterait a etablir.

Signalons que meme les analogues $l$-adiques des formules (4.6), (4.7), (4.8) ne sont pas prouves a l'heure actuelle; lorsque $y$ est la classe de cohomologie d'un cycle algebrique, (4.6) et (4.8) sont vraies en tous cas mod.\ torsion.


\subsubsection*{4.5.}
Dans cette section et les suivantes, nous allons donner des exemples ou dans le diagramme (4.1) l'homomorphisme $j$ respectivement l'homomorphisme $\phi$ n'est pas un isomorphisme. Il revient evidemment au meme, pour $X$ donne, de dire que
\[
j : \operatorname{Gr}^{\bullet}(X) \longrightarrow \operatorname{Gr}_{\mathrm{top}}^{\bullet}(X)
\]
est un isomorphisme, ou un monomorphisme, ou un epimorphisme. D'ailleurs l'image de $j$ est le sous-anneau de $\operatorname{Gr}_{\mathrm{top}}^{\bullet}(X)$ engendre par les classes de Chern des faisceaux localement libres sur $X$. Dans le cas ou $K^{\bullet}(X)$ est engendre comme algebre sur $\mathbb{Z}$ par les classes des faisceaux inversibles sur $X$, on voit tout de suite que le sous-anneau envisage de $\operatorname{Gr}_{\mathrm{top}}^{\bullet}(X)$ est aussi le sous-anneau engendre par $\operatorname{Gr}_{\mathrm{top}}^{1}(X)$, la composante de degre 1 de $\operatorname{Gr}_{\mathrm{top}}^{\bullet}(X)$.

Or prenons pour $X$ une variete de drapeaux
\[
X = G/B,
\]
ou $G$ est un groupe algebrique semi-simple simplement connexe sur le corps algebriquement clos $k$, et $B$ un sous-groupe de Borel [5]. Lorsque $k = \CC$, alors la decomposition cellulaire de Bruhat de $X$ [5, Exp.\ 13] implique que l'homomorphisme canonique
\[
K^{\bullet}(X) \longrightarrow K^{\bullet}(X(\CC)) \tag{$*$}
\]
est un isomorphisme, de meme d'ailleurs que l'homomorphisme canonique
\[
A(X) = H^{2\bullet}(X(\CC), \mathbb{Z}) \longrightarrow H^{2\bullet}(X(\CC), \mathbb{Z}),
\]
ou $X(\CC)$ est l'espace compact des points complexes de $X$, muni de sa topologie habituelle deduite de la topologie de $\CC$. D'autre part, il a ete prouve par L.\ HODGKIN [12] que $K^{\bullet}(X(\CC))$ est bien engendre par les classes des faisceaux inversibles sur $X$. Il en est donc de meme de $K^{\bullet}(X)$ lui-meme, et il n'est d'ailleurs pas difficile d'en deduire, par une technique de specialisation utilisant les groupes semi-simples deployes sur $\mathbb{Z}$ (SGA 3), que le meme resultat reste valable pour tout corps $k$ algebriquement clos. D'ailleurs, la decomposition cellulaire de $X$ (ou la bijectivite de $(*)$ dans le cas $k = \CC$) implique aisement que $A(X)$ est sans torsion, donc grace a 4.2 l'homomorphisme $\phi : A(X) \to \operatorname{Gr}_{\mathrm{top}}^{\bullet}(X)$ est un isomorphisme. De ces faits il resulte que $j$ est un isomorphisme si et seulement si $A(X)$ est engendre par ses elements de degre 1, ou encore, si $k = \CC$, si et seulement si la cohomologie entiere de $X(\CC)$ est engendree par les elements de degre 2, i.e.\ si le classifiant $BG$ est sans torsion. Il est connu d'autre part [15, p.\ 217 th.\ B (1)$\Leftrightarrow$(2), et 2.5] que ceci est le cas si et seulement si $G$ est ``sans torsion'', ou encore si $G$ est isomorphe a un produit de groupes $Sl(n)$ et $Sp(m)$. Il suffit donc de prendre pour $G$ un groupe $SO(n)$ avec $n \geq 7$, ou un groupe simple exceptionnel, pour avoir un exemple ou $j$ n'est pas un isomorphisme.

\subsubsection*{4.6.}
On notera que dans cet exemple, il est hors de doute que le ``bon'' anneau est l'anneau sans torsion $A(X) = \operatorname{Gr}_{\mathrm{top}}^{\bullet}(X)$, ayant comme base les cellules de la decomposition de Bruhat, anneau egalement isomorphe a la cohomologie entiere de $X(\CC)$ si $k = \CC$; par contre, l'anneau $\operatorname{Gr}^{\bullet}(X)$ apparait comme un anneau ayant de la torsion de nature assez pathologique. Cet exemple confirme donc la conception suivant laquelle l'invariant $\operatorname{Gr}^{\bullet}(X)$ n'est guere interessant que modulo torsion, i.e.\ tensorise par $\mathbb{Q}$. On voit en meme temps, a l'aide de cet exemple, que ce n'est que tensorise par $\mathbb{Q}$ que $\operatorname{Gr}^{\bullet}(X)$ a un caractere covariant par rapport aux morphismes propres, meme en se bornant aux immersions fermees de schemas quasi-projectifs et lisses. En effet, si $X$ est tel que $j$ ne soit pas un isomorphisme, il existe dans $X$ un sous-schema ferme integre $Y$ tel que $c_i^{\gamma}(Y) \in K^{\bullet}(X)$ ne soit pas dans $\Fil_{\mathrm{top}}^{i}(X)$, ou $i = \codim(Y, X)$. Nous pouvons prendre $Y$ de dimension minimale, de sorte que $\Fil_{\mathrm{top}}^{i+1}(X) = \Fil_{\mathrm{top}}^{i+2}(X) = \dots$. Si $Z$ est l'ensemble singulier de $Y$, et $X' = X - Z$, $Y' = Y - Z$, on en conclut que $c_i^{\gamma}(Y') \in K^{\bullet}(X')$ n'est pas dans $\Fil_{\mathrm{top}}^{i}(X')$. Compte tenu que $K^{\bullet}(X) \to K^{\bullet}(X')$ est surjectif et compatible avec les augmentations, et que son noyau est $\Fil_Z(K_{\bullet}(X)) \supset \Fil_{\mathrm{top}}^{\codim(Z,X)}(X) \supset \Fil_{\mathrm{top}}^{i+1}(X)$. Considerons alors l'inclusion
\[
f : Y' \hookrightarrow X',
\]
qui est une immersion fermee de codimension $i$ de schemas quasi-projectifs et lisses sur $k$, telle que
\[
f_{*} : K^{\bullet}(Y') \longrightarrow K^{\bullet}(X')
\]
ne soit pas compatible avec les $\gamma$-filtrations (modulo decalage par $i$), puisque $f_{*}(1)$ n'est pas dans $\Fil_{\mathrm{top}}^{i}(X')$.

\subsubsection*{4.7.}
Gardant les notations de 4.5, notons maintenant qu'il resulte de (Exp.\ IX) que le theoreme de HODGKIN cite equivaut au fait que l'on a un isomorphisme
\[
K^{\bullet}(G) = \mathbb{Z},
\]
donc $\operatorname{Gr}^{\bullet}(G) = \mathbb{Z}$, l'homomorphisme $K^{\bullet}(G/B) \to K^{\bullet}(G)$ etant a priori surjectif et ayant comme noyau l'ideal engendre par les $c_1(L) - 1$, ou $L$ parcourt les Modules inversibles sur $X = G/B$ associes aux representations $B \to \GG_m$ (qui donnent en fait tous les modules inversibles sur $X$ [5, Exp.\ 15]). On voit d'autre part par un raisonnement analogue que
\[
A(G/B) \longrightarrow A(G)
\]
est egalement surjectif, et que son noyau est l'ideal engendre par les elements de $A^{1}(G/B)$ definis par les representations $B \to \GG_m$, donc en fait par $A^{1}(G/B)$ tout entier; cela montre que $A(G)$ est reduit a $\mathbb{Z}$ si et seulement si ``$G$ est sans torsion'', i.e.\ $G$ est isomorphe a un produit de groupes $Sl(n)$ et $Sp(m)$. Donc dans ce cas, et dans ce cas seulement, l'homomorphisme
\[
\phi : A(G) \longrightarrow \operatorname{Gr}_{\mathrm{top}}^{\bullet}(G)
\]
est un isomorphisme. Donc si $G$ a de la torsion, le schema sous-jacent fournit un exemple d'un $\phi$ non bijectif.

\subsubsection*{4.8.}
On trouve de meme un exemple ou l'homomorphisme $\psi$ de (4.1)
\[
\psi : A(X) \longrightarrow H^{2\bullet}(X, \mathbb{Z}_l(\bullet)) \tag{4.9}
\]
ne se factorise pas par $\operatorname{Gr}_{\mathrm{top}}^{\bullet}(X)$, ou ce qui revient au meme ici, un exemple ou cet homomorphisme n'est pas nul en degres $> 0$. Se limitant (comme il est loisible grace a des theoremes de specialisation convenables) au cas ou $k = \CC$, on est ramene a etudier l'homomorphisme
\[
A(G) \longrightarrow H^{2\bullet}(G(\CC), \mathbb{Z}), \tag{4.10}
\]
deduit d'ailleurs par passage au quotient de l'homomorphisme canonique
\[
A(G/B) = H^{2\bullet}(G/B(\CC), \mathbb{Z}) \longrightarrow H^{2\bullet}(G(\CC), \mathbb{Z}).
\]
Il est tentant de conjecturer que l'homomorphisme (4.10) est en fait injectif, puisqu'on sait par ailleurs que les nombres premiers de torsion pour $G(\CC)$ sont exactement ceux qui divisent l'ordre du groupe fini
\[
\pi_1(G) = P/Q,
\]
o\`u $P$ est le poids et $Q$ le reseau radiciel. S'il en etait ainsi, des que $G$ aurait de la torsion, l'homomorphisme (4.10) ne se factoriserait pas par $\phi$, donc pour $l$ convenable (de facon precise, pour $l$ divisant l'ordre de $A^{2}(G)$, i.e.\ $l$ coefficient de torsion de $G$) il en serait de meme pour (4.9). On peut du moins affirmer que dans le cas du groupe $G = G_2$ sur le corps $\CC$, l'homomorphisme (4.10) n'est pas nul en degre $> 0$, de facon precise il existe un element de $A^{3}(G)$ dont l'image dans $H^{6}(G(\CC), \mathbb{Z})$ est d'ordre 2; il en resulte donc que l'homomorphisme (4.9) ne se factorise pas par $\phi$ lorsque $l = 2$. Le fait utilise concernant $G_2$, qui revient a dire qu'il existe dans $H^{6}(G/B(\CC), \mathbb{Z})$ un element dont l'image dans $H^{6}(G(\CC), \mathbb{Z})$ est non nulle, se trouve dans [3], modulo le dictionnaire connu entre le langage groupes algebriques semi-simples complexes, et groupes de Lie compacts; il m'a ete communique obligeamment par J.P.\ SERRE.

\subsection*{5. Complements sur les homomorphismes de Chern}
\addcontentsline{toc}{subsection}{5. Complements sur les homomorphismes de Chern}

\subsubsection*{5.1.}
L'exemple donne dans 4.5 n'exclut pas qu'on puisse peut-etre trouver, pour un schema $X$ quelconque, un homomorphisme d'anneaux
\[
g^{*} : \operatorname{Gr}^{\bullet}(X) \longrightarrow H^{2\bullet}(X, \mathbb{Z}_l(\bullet)),
\]
tels que pour tout Module localement libre $E$ sur $X$, $g^{*}$ transforme la classe de Chern $c_i$ de Exp.\ V en la classe de Chern cohomologique de SGA 5 VII. Bien entendu, ici $l$ designe un nombre premier distinct des caracteristiques residuelles de $X$. Comme l'anneau $\operatorname{Gr}^{\bullet}(X)$ est engendre par les $c_i(E)$, on voit en tous cas que $g^{*}$ est uniquement determine par la condition d'etre compatible avec les classes de Chern.

\subsubsection*{5.2.}
Utilisant l'homomorphisme de Chern
\[
K^{\bullet}(X) \longrightarrow H^{2\bullet}(X, \mathbb{Z}_l(\bullet)),
\]
on en conclut, par passage aux gradues associes respectivement a la filtration du premier membre et la filtration par le degre du second, un homomorphisme de groupes gradues (mais pas d'anneaux !)
\[
\gamma^{*} : \operatorname{Gr}_{\mathrm{top}}^{\bullet}(X) \longrightarrow H^{2\bullet}(X, \mathbb{Z}_l(\bullet)),
\]
qui satisfait a la relation
\[
\gamma^{*}(\cl(E)) = c_i^{(l)}(E) \mod \text{ torsion},
\]
ou $c_i^{(l)}(E)$ designe la classe de Chern $l$-adique. On en conclut, par division a partir des $\gamma^{*}$, un homomorphisme d'anneaux :
\[
g^{*} : \operatorname{Gr}^{\bullet}(X) \longrightarrow H^{2\bullet}(X, \mathbb{Z}_l(\bullet)) \otimes \mathbb{Z}[1/(i-1)!]
\]
transformant classe de Chern en classes de Chern; et tel que l'image du premier membre soit contenue dans celle de $H^{2\bullet}(X, \mathbb{Z}_l(\bullet))$; par suite, la question de l'existence de $g^{*}$ se resoud affirmativement pourvu qu'on remplace $H^{2\bullet}(X, \mathbb{Z}_l(\bullet))$ par son quotient par son sous-groupe de torsion. Le premier cas douteux pour definir les composantes $g^{i}$ de l'hypothetique homomorphisme $g$ est celui de $g^{2}$, ou n'arrive a priori qu'a definir $2c_2^{(l)} = \gamma^{*}(\cl(\Lambda^2 E - \Sym^2 E))$.

La reponse a la question d'existence de $g^{2}$ serait negative si par exemple on pouvait trouver trois faisceaux inversibles $L_1, L_2, L_3$ sur $X$ satisfaisant la relation
\[
(L_1 - 1)(L_2 - 1)(L_3 - 1) = 0 \text{ dans } K^{\bullet}(X),
\]
mais tels qu'on ait
\[
c_1^{(l)}(L_1) c_1^{(l)}(L_2) c_1^{(l)}(L_3) \neq 0 \text{ dans } H^{6}(X, \mathbb{Z}_l(3)).
\]

\subsubsection*{5.3.}
Lorsque $X$ est lisse et quasi-projectif sur un corps $k$, on peut de meme se poser la question de l'existence d'un homomorphisme d'anneaux
\[
g : \operatorname{Gr}^{\bullet}(X) \longrightarrow A(X)
\]
a valeurs dans l'anneau de Chow $A(X)$, transformant classes de Chern en classes de Chern, et on peut repeter a ce propos les reflexions qui precedent, qui montrent en particulier comment construire directement $(-1)^{i-1}(i-1)! \, g^{i}$ a defaut de $g^{i}$ lui-meme.

\subsubsection*{5.4.}
Signalons, pour terminer ce numero, qu'essentiellement la meme question se pose si, au lieu d'un schema $X$, on considere un espace topologique compact (disons) $X$, ce qui permet encore de considerer le gradue $\operatorname{Gr}^{\bullet}(X)$ de $K(X)$ muni de sa $\lambda$-filtration, et de se proposer de trouver un homomorphisme d'anneaux
\[
\operatorname{Gr}^{\bullet}(X) \longrightarrow H^{2\bullet}(X, \mathbb{Z})
\]
transformant classes de Chern en classes de Chern.


\subsection*{6. Theoreme de Riemann-Roch cohomologique, et homomorphisme de Gysin}
\addcontentsline{toc}{subsection}{6. Theoreme de Riemann-Roch cohomologique, et homomorphisme de Gysin}

\subsubsection*{6.1.}
Pour simplifier, nous nous bornerons ici aux schemas $X$ admettant un Module inversible ample, ce qui implique que l'on n'a pas a distinguer entre le $K^{\bullet}(X)$ naif et l'autre. Commencons par la formule de Riemann-Roch ``avec denominateurs'' pour un morphisme propre $f : X \to Y$ d'intersection complete. Comme (pour un nombre premier $l$ premier aux caracteristiques residuelles de $Y$) on peut definir les classes de Chern $l$-adiques d'un Module localement libre $E$,
\[
c_i^{(l)}(E) \in H^{2i}(X, \mathbb{Z}_l(i))
\]
(SGA 5 VII), correspondant a un homomorphisme
\[
K^{\bullet}(X) \longrightarrow H^{2\bullet}(X, \mathbb{Z}_l(\bullet)), \tag{6.1}
\]
la question se pose de donner une variante cohomologique de la formule de Riemann-Roch. Le seul ingredient qui manque pour donner un sens a la formule est un homomorphisme de Gysin
\[
f_{*} : H^{\bullet-2d}(X, \mathbb{Z}_l(\bullet-d)) \longrightarrow H^{\bullet}(Y, \mathbb{Z}_l(\bullet)) \tag{6.2}
\]
ou $d$ est la dimension relative virtuelle de $f$. Cet homomorphisme etant suppose defini, et tenant compte de la formule de Riemann-Roch sous la forme de Exp.\ VIII, a valeurs dans $\operatorname{Gr}^{\bullet}(Y)$, on constate que la validite de la formule de Riemann-Roch $l$-adique pour le morphisme $f$ equivaut exactement a la commutativite du carre
\[
\xymatrix{
\operatorname{Gr}^{\bullet}(X) \ar[r] \ar[d]_{f_{*}} & H^{2\bullet}(X, \mathbb{Q}_l(\bullet)) \ar[d]^{f_{*}} \\
\operatorname{Gr}^{\bullet-d}(Y) \ar[r] & H^{2(\bullet-d)}(Y, \mathbb{Q}_l(\bullet-d))
} \tag{6.3}
\]
ou les fleches horizontales se definissent a partir de (6.1) en passant aux gradues associes et en divisant, de facon a transformer classes de Chern en classes de Chern, et ou les fleches verticales sont respectivement la fleche de Exp.\ VIII et l'homomorphisme de Gysin en cohomologie (suppose donne).

\subsubsection*{6.2.}
Notons que l'homomorphisme de Gysin est defini en tous cas si $Y$ est lisse sur un corps $k$ (SGA 5 IV), en utilisant le theoreme de purete cohomologique relatif sur $k$; la meme construction s'applique lorsque $Y$ est excellent regulier et de caracteristique nulle, grace au theoreme de purete cohomologique absolu de SGA 4 XIX, et elle s'appliquera automatiquement pour tout $Y$ excellent et regulier, sans restriction de caracteristique, une fois le theoreme de purete cohomologique absolu prouve en toute generalite. Enfin, la meme construction s'applique, grace au theoreme de purete cohomologique relatif (SGA 4 XVI 3.7), lorsque $X$ et $Y$ sont tous deux lisses sur un meme schema $S$ et que $f$ est un $S$-morphisme. Il y a lieu de conjecturer que dans chacun de ces cas ou l'homomorphisme de Gysin est defini, le diagramme (6.3) est bien commutatif. Le seul cas ou cette assertion soit actuellement verifiee est celui ou $X$ et $Y$ sont tous les deux lisses sur un corps $k$, et resulte alors de la commutativite du diagramme analogue a (6.3) :
\[
\xymatrix{
A^{\bullet}(X) \ar[r] \ar[d]_{f_{*}} & H^{2\bullet}(X, \mathbb{Z}_l(\bullet)) \ar[d]^{f_{*}} \\
A^{\bullet-d}(Y) \ar[r] & H^{2(\bullet-d)}(Y, \mathbb{Z}_l(\bullet-d))
} \tag{6.4}
\]
faisant intervenir les anneaux de classes de cycles (mod.\ equivalence rationelle) $A(X)$ et $A(Y)$, commutativite etablie dans SGA 5 IV.

\subsubsection*{6.3.}
Notons que dans chacun des cas envisages plus haut ou on arrive a definir un homomorphisme de Gysin, on obtient en fait une construction pour la cohomologie a coefficients dans $\mathbb{Z}_l(i)$, et non seulement dans $\mathbb{Q}_l(i)$. Cela permet, lorsque $f$ est une immersion, de poser la question de la validite de la formule de Riemann-Roch ``sans denominateurs'', comme il est explique au n\textsuperscript{o} 3. Lorsque $X$ et $Y$ sont tous deux lisses sur un corps $k$, la commutativite de (6.4) nous montre qu'il suffirait de prouver la formule sans denominateurs dans l'anneau de Chow $A(Y)$ (cf.\ (3.3)).

\subsubsection*{6.4.}
Lorsqu'on ne fait pas d'hypothese de regularite sur $Y$, et meme si $X$ et $Y$ sont de type fini sur le corps $\CC$, on ne voit plus de definition plausible d'un homomorphisme de Gysin cohomologique (6.2), meme pour une immersion reguliere de codimension 1. Peut-on en trouver un neanmoins avec des proprietes raisonnables, en particulier rendant commutatif le carre (6.3) ? La question est liee evidemment a celle de definir, pour un $X$ d'intersection complete de codimension $i$ dans $Y$, ``sa classe de cohomologie''
\[
\cl(X) \in H^{2i}(Y, \mathbb{Z}_l(i)),
\]
de facon a transformer intersections (sans composantes excedentaires) en cup-produits. Cette derniere question se resoud affirmativement tout au moins modulo torsion, i.e.\ dans $H^{2i}(Y, \mathbb{Q}_l(i))$, en definissant
\[
\cl(X) = c_i(\cO_X) / (-1)^{i-1}(i-1)! \in H^{2i}(Y, \mathbb{Q}_l(i)).
\]
Mais j'ignore si $c_i(\cO_X)$ est divisible par $(i-1)!$ dans $H^{2i}(Y, \mathbb{Z}_l(i))$. En particulier, si $Y$ est un schema projectif de dimension $3$ sur le corps des complexes, disons, et $X$ reduit a un point (singulier) de $Y$, $c_3(\cO_X)$ est-il divisible par $2$ dans $H^{6}(Y, \mathbb{Z}_l(3))$, et en particulier la classe de Chern est-elle nulle dans $H^{6}(Y, \mathbb{Z}/2\mathbb{Z})$ ? Signalons egalement dans la meme voie la question suivante : soit $Y$ un schema projectif (singulier) de dimension $4$, $X$ un sous-schema ferme de dimension $1$ qui soit globalement une intersection de trois diviseurs positifs $D_1, D_2, D_3$, considerons le produit de leurs classes de Chern mod.\ $2$ :
\[
c_1(D_1) c_1(D_2) c_1(D_3) \in H^{6}(Y, \mathbb{Z}/2\mathbb{Z});
\]
cet element du n\textsuperscript{o} est-il independant du choix particulier des $D_i$ d'intersection $X$ ?

\subsection*{7. Classes de Chern des complexes parfaits}
\addcontentsline{toc}{subsection}{7. Classes de Chern des complexes parfaits}

\subsubsection*{7.1.}
Lorsqu'on se pose la question de developper une formule de Riemann-Roch cohomologique, avec ou sans denominateurs, en dehors de l'hypothese assez artificielle d'existence d'un Module inversible ample sur les schemas envisages, on est necessairement confronte a la question d'une definition, pour un complexe de Modules parfait $L_{\bullet}$ sur un schema $X$, des classes de Chern
\[
c_i(L_{\bullet}) \in H^{2i}(X, \mathbb{Z}_l(i)),
\]
generalisant les classes de Chern des Modules localement libres SGA 5 VII. La meme question se pose d'ailleurs pour des topos localement anneles generaux (comparer [10]), ainsi que pour les complexes parfaits sur des espaces paracompacts anneles par les fonctions complexes continues, en travaillant alors avec la cohomologie entiere ordinaire. Dans le cas d'un espace compact, on sait que le $K^{\bullet}(X)$ construit avec les complexes parfaits est le meme que le ``naif'' construit avec les Modules localement libres (Exp.\ IV), et dans ce cas on est ramene trivialement a la definition classique. Mais deja dans le cas d'un espace localement compact le probleme pose semble non trivial, puisque la cohomologie entiere d'un tel espace n'est pas toujours la limite projective des cohomologies de ses parties compactes.

\subsubsection*{7.2.}
Plus precisement, lorsque $Y$ est une partie fermee de $X$ telle que $L_{\bullet}|(X-Y)$ soit nul (au sens des categories derivees), i.e.\ telle que les faisceaux de cohomologie de $L_{\bullet}$ soient a supports dans $Y$, on aimerait definir des classes de Chern a support dans $Y$
\[
c_i^{Y}(L_{\bullet}) \in H_Y^{2i}(X, \mathbb{Z}_l(i)) \quad \text{resp.}\ H^{2i}(X, \mathbb{Z}_l(i))
\]
qui redonneraient les classes de Chern ``ordinaires'' de 7.1 a l'aide de l'homomorphisme canonique
\[
H_Y^{2i}(X, \mathbb{Z}_l(i)) \longrightarrow H^{2i}(X, \mathbb{Z}_l(i)).
\]
Meme dans le cas ou ces dernieres se ramenent aux classes de Chern des Modules localement libres ($X$ schema avec Module inversible ample, ou espace compact), la construction de ces classes de Chern ``a supports'' n'est pas faite, et ne pourra sans doute se faire qu'en meme temps que celle des classes de Chern generales postulees dans 7.1\footnote{Ajoute en octobre 1968. Les problemes souleves ici ont ete resolus dans le cadre de la cohomologie de Hodge par L.\ Illusie; cf.\ [13].}

Il est probable que les classes de Chern avec supports auront un role a jouer dans la demonstration de la formule de Riemann-Roch cohomologique sans denominateurs (cf.\ 3.5). D'autre part, elles devraient permettre de donner une version amelioree de la formule de Riemann-Roch (avec ou sans denominateurs) en introduisant dans cette formule des familles de supports.

\subsection*{8. Anneau de Chow des schemas reguliers}
\addcontentsline{toc}{subsection}{8. Anneau de Chow des schemas reguliers}

Le probleme est de definir l'equivalence rationelle des cycles, et d'etudier l'anneau de classes obtenu (notamment ses proprietes fonctorielles), pour un schema noetherien regulier $X$ admettant un Module inversible ample, en generalisant la theorie connue lorsque $X$ est lisse sur un corps [4]. Le resultat technique clef sera encore un ``moving lemma'', assurant qu'on peut faire varier un cycle dans sa classe de facon a intersecter avec la bonne dimension un cycle donne.

Le probleme envisage est assez naturel a priori, d'autant plus que nous avons vu (4.5 a 4.8) que les substituts $\operatorname{Gr}^{\bullet}(X)$, $\operatorname{Gr}_{\mathrm{top}}^{\bullet}(X)$ de l'anneau de Chow $A(X)$ ne sont guere satisfaisants quand on ne se borne pas a une theorie ``modulo torsion''. Lorsque $X$ se deduit d'un schema algebrique projectif et lisse $Y$ sur un corps $k$, en prenant le schema local au sommet d'un cone projetant $C$ et enlevant l'origine, on trouve (en admettant que $A(X)$ est bien la limite inductive des anneaux de Chow des $U - \{0\}$, ou $U$ parcourt les voisinages ouverts dans $C$ du sommet $O$) que $A(X)$ est canoniquement isomorphe a l'anneau quotient $A(Y)/\xi A(Y)$, ou $\xi \in A^{1}(Y)$ est induit par l'immersion projective consideree de $Y$. Cela montre que la connaissance de l'anneau de Chow $A(X)$ (de type non classique) equivaut a la connaissance de $A(Y)/\xi A(Y)$, qu'on pourrait appeler la ``partie primitive'' de l'anneau de Chow $A(Y)$ (en un sens suggere par la theorie cohomologique de Lefschetz-Hodge des varietes algebriques projectives non singulieres). Cet exemple montre que les ``anneaux de Chow locaux'' des anneaux locaux a singularite isolee sont des invariants geometriques tout a fait non triviaux, et qui auront sans doute leur role a jouer dans l'etude ``arithmetiques'' de ces anneaux locaux.

Bien entendu, une fois en possession d'une notion raisonnable d'anneau de Chow pour un schema noetherien regulier admettant un Module inversible ample, il y a lieu de generaliser a ce cas les relations etablies resp.\ conjecturees au n\textsuperscript{o} 4 entre cet anneau et l'anneau $\operatorname{Gr}_{\mathrm{top}}^{\bullet}(X)$, et a formuler en termes de cet anneau la question de la ``forme de Riemann-Roch sans denominateurs'' enoncee au n\textsuperscript{o} 3.

\subsection*{Bibliographie}
\addcontentsline{toc}{subsection}{Bibliographie}

\begin{enumerate}[label={[\arabic*]}]
\item[{[1]}] M.F.\ Atiyah et F.\ Hirzebruch, The Riemann-Roch theorem for analytic embeddings, Topology, vol.\ 1, p.\ 151--166 (1962).
\item[{[2]}] A.\ Borel et J.P.\ Serre, Le theoreme de Riemann-Roch, Bull.\ Soc.\ Math.\ France, vol.\ 86, p.\ 97--136 (1958).
\item[{[3]}] R.\ Bott et H.\ Samelson, Applications of the theory of Morse to symmetric spaces, Amer.\ Journ.\ of Math.\ vol.\ 80, p.\ 964--1029 (1958).
\item[{[4]}] C.\ Chevalley, Les classes d'equivalence rationelle, I et II, Seminaire a l'ENS 1958, exposes 2 et 3.
\item[{[5]}] C.\ Chevalley, Classification des groupes de Lie algebriques, Seminaire a l'ENS 1956/58.
\item[{[6]}] A.\ Dold et D.\ Puppe, Homologie nicht-additiver Funktoren.\ Anwendungen, Annales de l'Inst.\ Fourier t.\ XI, p.\ 201--312 (1961).
\item[{[7]}] H.\ Grauert, Ein Theorem der analytischen Garbentheorie, Pub.\ Math.\ I.H.\'E.S.\ n\textsuperscript{o} 5, p.\ 5--64 (1960).
\item[{[8]}] A.\ Grothendieck, Classes de faisceaux et theoreme de Riemann-Roch, notes multigraphiees par les soins de l'Institute for Advanced Study, 1957 (travail cite [RRR] dans le present expose, et qui est reproduit dans notre seminaire sous forme d'Appendice au Chap.\ 0).
\item[{[9]}] A.\ Grothendieck, La theorie des classes de Chern, Bull.\ Soc.\ Math.\ France vol.\ 86, p.\ 137--154 (1958).
\item[{[10]}] A.\ Grothendieck, Classes de Chern et representations lineaires de groupes discrets, in : Dix exposes sur la cohomologie des schemas, Masson--North Holland Pub.\ Cie.
\item[{[11]}] H.\ Hironaka, Resolution of singularities of an algebraic variety over a field of characteristic zero, Ann.\ of Math.\ 79, p.\ 109--326 (1964).
\item[{[12]}] L.\ Hodgkin, (a paraitre).
\item[{[13]}] L.\ Illusie, travail en preparation.
\item[{[14]}] J.P.\ Serre, Algebre Locale, Multiplicites, Lecture Notes in Math.\ n\textsuperscript{o} 11 (1965) Springer.
\item[{[15]}] A.\ Borel, Sous-groupes commutatifs et torsion des groupes de Lie compacts connexes, Tohoku Math.\ Journal, vol.\ 13, p.\ 216--240 (1961).
\item[{[16]}] P.\ Deligne, Exp.\ XVIII in Theorie des topos et cohomologie etale des schemas (SGA 4), par M.\ Artin, A.\ Grothendieck et J.L.\ Verdier.
\item[{[17]}] Yu.\ Manin, Correspondances, motifs et transformations monoidales (en russe), Mat.\ Sbornik, t.\ 77 (119) n\textsuperscript{o} 4, 1968.
\end{enumerate}

\vspace{1cm}
\noindent\textbf{Note pour la page 1.} La theorie des operations $\lambda$ dans $D(X)$, et des operations $\sigma_i$ dans $K^{\bullet}(X)$, a ete ecrite par P.\ Deligne en 1968. Il prouve que $K^{\bullet}(X)$ devient ainsi un pre-$\lambda$-anneau, et lorsque $\cO_X$ est un faisceau d'Algebres sur un corps, que $K^{\bullet}(X)$ est meme un $\lambda$-anneau. Il est plausible qu'un raffinement de la methode de Deligne permette de se debarrasser de l'hypothese restrictive sur $\cO_X$.

\vspace{0.5cm}
\noindent\textbf{Notes pour la page 11}

\noindent $(x)$ (Avril 1969) D.\ MUMFORD nous signale une demonstration de (4.6) dans l'anneau de Chow (sans negliger la torsion), qu'il a obtenue dans un seminaire a Harvard en 1959 ! Le lecteur trouvera la demonstration de ce beau resultat dans SGA 5 VII 9. Ce resultat contient evidemment le resultat cohomologique correspondant, signale dans la note $(x*)$. Cependant, le resultat cohomologique est prouve par J.P.\ JOUANOLOU sous des conditions plus generales (sans hyperplan quasi-projectif, et dans le cas lisse relatif sur une base quelconque).

\vspace{0.3cm}
\noindent $(xx)$ (Janvier 1969) J.P.\ JOUANOLOU vient de prouver ces formules pour tout $y$ sans negliger la torsion, cf.\ SGA 5 VII \S\ 4.9. DELIGNE et MANIN [17] avaient auparavant remarque independamment, et par des arguments (faciles) de nature assez differente, qu'il suffisait d'etablir des formules pour $y$ ``algebrique'', et meme pour $y = 1$ dans (4.6), pour conclure la meme relation dans le cas general.


% ============================================================
\newpage
\section*{Index terminologique}
\addcontentsline{toc}{section}{Index terminologique}

\begin{longtable}{p{8cm}l}
\toprule
\textbf{Terme} & \textbf{Expos\'e} \\
\midrule
\endhead

additive (application) & I \\
additive (application) & III \\
$\lambda$-alg\`ebre augment\'ee & V \\
alg\'ebriquement \'equivalent \`a z\'ero (dans $K^{\bullet}$) & XIII \\
alg\'ebriquement \'equivalent \`a z\'ero (faisceau inversible) & XIII \\
d'amplitude $\mathfrak{C}$-parfaite contenue dans $[a,b]$ (complexe) & I \\
d'amplitude $\mathfrak{C}$-parfaite finie (complexe) & I \\
d'amplitude $\mathfrak{C}$-parfaite contenue dans $[a,b]$ (complexe) & III \\
d'amplitude parfaite finie (complexe) & I \\
d'amplitude plate contenue dans $[a,b]$ (complexe) & I \\
d'amplitude plate finie (complexe) & I \\
d'amplitude parfaite relativement \`a $f$ (ou $Y$) contenue dans $[a,b]$ (complexe) & III \\
d'amplitude parfaite contenue dans $[-n,0]$ (morphisme) & III \\
d'amplitude parfaite finie relativement \`a $f$ (complexe) & III \\
d'amplitude plate relativement \`a $f$ (ou $Y$) contenue dans $[a,b]$ (complexe) & III \\
d'amplitude plate contenue dans $[-n,0]$ (morphisme) & III \\
d'amplitude plate finie relativement \`a $f$ (complexe) & III \\
amplitude projective & II App.\ II \\
$\lambda$-anneau & V \\
$\lambda$-anneau libre engendr\'e par une famille de g\'en\'erateurs & V \\
anneau de Chern & V \\
application polyn\^ome & X \\
bin\^omial (anneau) & V \\
$\lambda$-anneau engendr\'e par g\'en\'erateurs et relations & V \\
caract\`ere de Chern & V \\
caract\`ere de Todd & V \\
$S$-cat\'egorie ab\'elienne & I \\
$S$-cat\'egorie ab\'elienne plate & I \\
$S$-cat\'egorie additive & I \\
$S$-cat\'egorie triangul\'ee & I \\
classe de Chern & V \\
classe de Chern compl\'et\'ee & V \\
classe de Chern totale & V \\
codimension (d'une immersion r\'eguli\`ere) & VII \\
coh\'erent (complexe) & I \\
coh\'ereur & II \\
complexe cotangent relatif & VIII \\
complexe elliptique & II App.\ II \\
cup-produit & IV \\
$k$-cycle positif sp\'ecial & XIII \\
degr\'e (d'une application polyn\^ome) & X \\
degr\'e (d'un faisceau inversible) & XIII \\
degr\'e (d'un morphisme de sch\'emas) & X \\
dense (sous-ensemble d'une cat\'egorie ab\'elienne ou triangul\'ee) & IV \\
dimension parfaite & I \\
dimension relative virtuelle & VIII \\
divisoriel & II \\
dual & I \\
dualit\'e parfaite & I \\
\'equivalence num\'erique (dans $\operatorname{Gr}^{\bullet}(X)$ et $\operatorname{Gr}_{\bullet}(X)$) & XIII \\
$\tau$-\'equivalence & XIII \\
espace pseudo-analytique complexe & II App.\ II \\
espace pseudo-analytique complexe de Stein & II App.\ II \\
exact (sous-ensemble d'une cat\'egorie ab\'elienne ou triangul\'ee) & IV \\
$\mathrm{(b)}$-faisceau & XIII \\
faisceau conormal & VII \\
faisceau cotangent relatif & II App.\ II \\
faisceau tangent relatif & II App.\ II \\
$\mathrm{(b)}$-famille & XIII \\
famille ample & XIII \\
famille de classes de faisceaux coh\'erents & XIII \\
famille de classes de faisceaux coh\'erents limit\'ee & XIII \\
famille de classes de faisceaux coh\'erents limit\'ee par un faisceau coh\'erent & XIII \\
fibr\'e cotangent relatif & II App.\ II \\
fibr\'e tangent relatif & II App.\ II \\
fid\`element plat (morphisme de topos annel\'es) & I \\
filtration de $\lambda$-alg\`ebre augment\'ee & V \\
$S$-foncteur additif & I \\
$S$-foncteur exact & I \\
formule de projection & IV \\
groupe de Grothendieck & IV \\
$\lambda$-homomorphisme & V \\
homomorphisme de sp\'ecialisation & I \\
$\lambda$-id\'eal & V \\
indice analytique & II App.\ II \\
d'intersection compl\`ete (morphisme) & VIII \\
$n$-isomorphisme & I \\
libre homog\`ene (module gradu\'e) & I \\
localement stable par noyau d'\'epimorphisme & I \\
localement quasi-relevable & I \\
localement quasi-stable par noyau d'\'epimorphisme & I \\
localement relevable & I \\
nombre de Chern & V \\
nombre d'intersection & XIII \\
nombre d'intersection (pour une famille de faisceaux inversibles) & XIII \\
num\'eriquement \'equivalent \`a z\'ero (faisceau inversible) & XIII \\
op\'erateur diff\'erentiel universel d'ordre $n$ & II App.\ II \\
op\'erations d'Adams & V \\
$\mathfrak{C}$-parfait (complexe) & I \\
parfait (complexe) & I \\
parfait (morphisme) & III \\
parfait relativement \`a $f$ (complexe) & III \\
pr\'ecoh\'erent & I \\
de $n$-$\mathfrak{C}$-pr\'esentation finie & I \\
de pr\'esentation finie & I \\
d'\'ec-$\mathfrak{C}$-pr\'esentation finie & I \\
$n$-$\mathfrak{C}$-pseudo-coh\'erent (complexe) & I \\
$\mathfrak{C}$-pseudo-coh\'erent (complexe) & I \\
$n$-pseudo-coh\'erent (complexe) & I \\
$n$-pseudo-coh\'erent (morphisme) & III \\
pseudo-coh\'erent (complexe) & I \\
pseudo-coh\'erent (morphisme) & III \\
$n$-pseudo-coh\'erent relativement \`a $f$ (complexe) & III \\
pseudo-coh\'erent relativement \`a $f$ (complexe) & III \\
$\mathrm{(b)}$-polyn\^ome & XIII \\
polyn\^ome de Snapper & X \\
polyn\^ome universel & V \\
pr\'e-$\lambda$-anneau & V \\
$n$-quasi-isomorphisme & I \\
rang (d'un complexe parfait) & X \\
$m$-r\'egulier (faisceau) & XIII \\
r\'egulier (homomorphisme) & VII \\
r\'egulier (id\'eal) & VII \\
r\'eguli\`ere (immersion) & VII \\
sous-$S$-cat\'egorie triangul\'ee & I \\
sp\'ecial ($\lambda$-anneau) & V \\
sp\'ecial (morphisme) & III \\
de Stein (topos localement annel\'e) & II App.\ II \\
strictement $\mathfrak{C}$-parfait & I \\
strictement $\mathfrak{C}$-$n$-pseudo-coh\'erent & I \\
strictement $\mathfrak{C}$-pseudo-coh\'erent & I \\
syst\`eme de g\'en\'erateurs & VII \\
\bottomrule
\end{longtable}


\newpage
\section*{Index des notations}
\addcontentsline{toc}{section}{Index des notations}

\begin{longtable}{p{7cm}l}
\toprule
\textbf{Notation} & \textbf{Expos\'e} \\
\midrule
\endhead

$\ch$ & V \\
$\Ch(A)$, $\Ch_{\bullet}(A)$ & I \\
$\Coh(S)$ & I \\
$\mathfrak{C}_n$-$\mathrm{parf\text{-}amp}(F)$ & II App.\ II \\
$c^{(l)}(E)$, $c_i^{(l)}(E)$ & II App.\ II \\
$c_i(E)$, $c_i(x)$ & V \\
$D^{\mathfrak{C}}(S)$, $D^{+}(S)$, $D^{-}(S)$, $D^{b}(S)$ & I \\
$D^{\mathfrak{C}}_{\mathrm{qc}}(S)$, $D_{\mathrm{qc}}(S)$, $D_{\mathrm{qc}, n\text{-}\mathrm{cons}}(S)$ & III \\
$D^{\mathfrak{C}}_{\mathrm{parf}}(S)$, $D^{\mathfrak{C}}_{\mathrm{parf, cons}}(S)$, $D_{\mathrm{parf}}(S)$ & III \\
$D^{\mathfrak{C}}_{n\text{-}\mathrm{pseudo\text{-}coh}}(S)$, $D^{\mathfrak{C}}_{\mathrm{pseudo\text{-}coh}}(S)$, $D_{n\text{-}\mathrm{pseudo\text{-}coh}}(S)$ & III \\
$D_{\mathrm{parf}}^{\mathfrak{C}}(X)$, $D_{\mathrm{parf}}^{\mathfrak{C}, f\text{-}\mathrm{parf}}(X)$, $D^{\mathfrak{C}}(X)$, $D(X)$ & I \\
$D_{\mathrm{parf}}^{\mathfrak{C}, f\text{-}\mathrm{torf}}(X)$, $D_{\mathrm{parf}}^{\mathfrak{C}}(X)$ & III \\
$D_{n\text{-}\mathrm{cons}}(X)$, $D_{\mathrm{cons}}(X)$, $D_{\mathrm{parf}}(X)$, $D_{\mathrm{torf}}(X)$ & III \\
$D_{\mathrm{coh}}^{\mathfrak{C}}(X)$, $D_{\mathrm{coh}}(X)$ & I \\
$E^{\vee}$ & I \\
$E \otimes F$, $E \otimes^{\LL} F$ & I \\
$f^{*}$ (entre les $K_{\bullet}$) & IV \\
$f^{*}$ (entre les $K^{\bullet}$) & IV \\
$f_{*}$ (entre les $K_{\bullet}$) & IV \\
$f_{!}$ (entre les $K^{\bullet}$) & IV \\
$f_{*}$ (cas du fibr\'e projectif) & VI \\
$f_{!}$ & IV \\
$\mathrm{Filt}_{\top}(X)$ & X \\
$\operatorname{Gr}^{\bullet}(X)$, $\operatorname{Gr}_{\bullet}(X)$ & V, X \\
$\operatorname{Gr}_{\mathrm{top}}^{\bullet}(X)$ & XIV \\
$\operatorname{Gr}^{i}(X)$ & V \\
$\Hom^{\bullet}(E, F)$ & I \\
$\Hom(E, F)$ & I \\
$K^{\bullet}(X)$, $K_{\bullet}(X)$ & IV \\
$K^{\bullet}(X)_{\mathrm{naif}}$, $K_{\bullet}(X)_{\mathrm{naif}}$ & IV \\
$K^{\bullet}(f)$, $K_{\bullet}(f)$ & IV \\
$K^{\bullet}_{\mathrm{num}}(X)$, $K^{\bullet}_{\mathrm{alg}}(X)$ & XIII \\
$K^{\bullet}(X/Y)$, $K_{\bullet}(X/Y)$ & IV \\
$\mathrm{Lib}(U)$, $\mathrm{Lib}_{\mathrm{hom}}(U)$ & I \\
$\Lambda^{i}(E)$, $\Lambda^{i}(L_{\bullet})$ & V, XIV \\
$\lambda^{i}(x)$, $\lambda_{t}(x)$ & V \\
$\lambda_{t}^{-1}(M)$ & V \\
$\lambda^{i}_{T}(x)$ & V \\
$\mathcal{M}\mathrm{od}(U)$ & I \\
$\mathrm{ob}\,\mathcal{C}$ & I \\
\bottomrule
\end{longtable}

\vspace{1cm}

\begin{longtable}{p{7cm}l}
\toprule
\textbf{Notation} & \textbf{Expos\'e} \\
\midrule
\endhead

$\mathrm{par\text{-}amp}_{\bullet}(E)$, $\mathrm{par\text{-}amp}^{\bullet}(E)$ & XIII \\
$\mathrm{Parf}(\mathfrak{C}, S)$, $\mathrm{Parf}(\mathcal{C}, \mathcal{C}')$, $\mathrm{Parf}(S)$ & I, III \\
$\Pic^{\tau}(X)$, $\Pic^{\tau}_{X/S}$ & XIII \\
$\Pic^{0}(X)$, $\Pic^{0}_{X/k}$ & XIII \\
$\Pic^{\mathrm{num}}(X)$, $\Pic^{\mathrm{num}}_{X/S}$ & XIII \\
$\Pic^{[n]}(X)$ & XIII \\
$\Pic(X)$, $\Pic_{X/S}$ & XIII \\
$P_{X/S}$ & XIII \\
$P_{X/S}^{(n)}$ & XIII \\
$P_{X/S}(E)$ & XIII \\
$\Proj$ & VI \\
$\mathcal{P}_{\gamma, \infty}(X)$, $\mathcal{P}_{\lambda, \infty}(X)$ & V \\
$\mathrm{Qcoh}(S)$ & I \\
$T^{\bullet}(E)$, $T_{\bullet}(E)$ & VIII \\
$\Todd(E)$, $\Todd(f)$, $\Td(E)$ & VIII \\
$\mathrm{tor\text{-}amp}(F)$, $\mathrm{tor\text{-}amp}_{\bullet}(E)$, $\mathrm{tor\text{-}amp}^{\bullet}(E)$ & II App.\ II \\
$\mathrm{Trip}$ & I \\
$X/S$ & III \\
$X^{(n)}$ & V \\
\bottomrule
\end{longtable}

\vspace{1cm}

\begin{longtable}{p{7cm}l}
\toprule
\textbf{Notation} & \textbf{Expos\'e} \\
\midrule
\endhead

$1 + A[[t]]^{*}$ & V \\
$Z(x)$ & V \\
\bottomrule
\end{longtable}

\vfill
\begin{center}
\rule{0.5\textwidth}{0.4pt}

\medskip
\textit{Fin du S\'eminaire de G\'eom\'etrie Alg\'ebrique du Bois Marie 1966--67}
\end{center}

\end{document}
